% concatenated from BHMP_Ki_Po.tex
\documentclass[reqno,11pt]{amsart}
\usepackage{amsthm,amsfonts,amssymb,euscript,mathrsfs,graphics,color,amsmath,latexsym,marginnote,hyperref}
%,todonotes,multicol,yhmath,imakeidx}
\usepackage[latin1]{inputenc}
%\usepackage{graphicx}
%\usepackage[english]{babel}
%\usepackage[hypertex]{hyperref} % buono per DVI 

%\usepackage[dvips]{graphicx} 
%\usepackage[margin=1in]{geometry}

%\usepackage{showkeys}
%\usepackage{showlabels}
\usepackage{amsmath}
\usepackage{graphicx}
\usepackage{amsfonts}
\usepackage{amssymb}
\usepackage{amsthm}
\usepackage{dsfont}    
\usepackage{mathrsfs}
\usepackage{amsthm}
\usepackage{verbatim}
\usepackage{esint}
\usepackage{stmaryrd}
\usepackage{tcolorbox}

\usepackage[color]{changebar}
\usepackage{marginnote}
\usepackage{xcolor}

\cbcolor{red!70!black}        % colore della barra
\setlength{\changebarwidth}{1.2pt}
\setlength{\changebarsep}{6mm} % distanza barra?testo

% \nota[colore]{testo della nota}{testo annotato}
\newcommand{\nota}[3][red!70!black]{%
\cbstart\marginnote{\footnotesize\color{#1}\raggedright #2}#3\cbend{}}


\theoremstyle{plain}

\def\tofill{\vskip30pt $\cdots$ To fill in $\cdots$ \vskip30pt}

\usepackage{hyperref}
%\usepackage[]{hyperref}
%\hypersetup{
%    colorlinks=false,       
%    linkcolor=red,         
%    citecolor=blue,}

%\usepackage{mathptmx}
\usepackage{times}
%\usepackage{lmodern} % math, rm, ss, tt 
%\usepackage[T1]{fontenc} 
%\usepackage{fourier}
%\usepackage{smfthm}
%\usepackage{makeidx}
%\usepackage{showidx}
\usepackage{todonotes}
\usepackage{mathtools}

\newcommand{\todog}[1]{\todo[color=red]{#1}}
\setlength{\marginparwidth}{3cm}
\usepackage{marginnote}

%\usepackage[notref,notcite]{showkeys}

\numberwithin{equation}{section}

%\textwidth 19cm
%\textheight 22.5cm %\oddsidemargin 0cm
%\evensidemargin 0cm
%\topmargin 0.05cm
%\headheight 0.38cm
%\headsep 0.85cm
%\footskip 2.35cm
%Valori standard:         art10          bk12       rep12

\oddsidemargin   0cm  %        2.21cm |       0.63 cm |  0.74 cm
\evensidemargin 0cm  %        2.21   |       3.17    |  2.07
\topmargin  0.05cm      %        0.95   |       1.85    |  0.95
\headheight 0.37cm      %        0.42   |       0.69    |  0.42
\headsep    0.38cm      %        0.88   |       0.69    |  0.88
\textwidth  16.5cm      %       12.12   |       12.7    | 13.7
\textheight 22.5cm     %       18.55   |       18.85   | 18.85
\footskip   1.5cm      %        1.05   |       1.05    |  1.05
%
%       deve essere  oddside + evenside + textwidth = 16.5cm



\newenvironment{blocco}[1][]{
\begin{tcolorbox}[
	colback=blue!15,        % colore sfondo principale
	colframe=blue!50,       % colore bordo
	colbacklower=blue!30,   % colore sfondo più scuro in basso
	title= \ifx&#1& \else#1\fi,
	fonttitle=\bfseries,
	arc=4pt,                % arrotondamento angoli
	boxrule=0.8pt           % spessore bordo
	]
}{
\end{tcolorbox}
}

\tcbuselibrary{skins}


\newenvironment{pf}{\noindent{\sc Proof}.\enspace}{\rule{2mm}{2mm}\medskip}
\newenvironment{pfn}{\noindent{\sc Proof} \enspace}{\rule{2mm}{2mm}\medskip}
\newtheorem{thm}{Theorem}[section]
\newtheorem{rmk}[thm]{Remark}
\newtheorem{prop}[thm]{Proposition}
\newtheorem{cor}[thm]{Corollary}
\newtheorem{theorem}{Theorem}[section]
\newtheorem{proposition}[theorem]{Proposition}
\newtheorem{lemma}[theorem]{Lemma}
\newtheorem{corollary}[theorem]{Corollary}
\newtheorem{condition}[theorem]{Condition}
\newtheorem{remark}[theorem]{Remark}
\newtheorem{remarks}[theorem]{Remark}
\newtheorem{hypo}[theorem]{Hypothesis}
\newtheorem{definition}[theorem]{Definition}
\newtheorem{defn}[thm]{Definition}
\newenvironment{Remark}{\begin{remark} \rm}{\rule{2mm}{2mm}\end{remark}}
\newenvironment{Remarks}{\begin{remarks} \rm}{\rule{2mm}{2mm}\end{remarks}}
\newcommand{\R}{\mathbb{R}}
\newcommand{\N}{\mathbb{N}}
\newcommand{\Z}{\mathbb{Z}}
\newcommand{\Q}{\mathbb{Q}}
\newcommand{\C}{\mathbb{C}}
\newcommand{\T}{\mathbb{T}}
\newcommand{\M}{\mathcal{M}}
\renewcommand{\L}{\mathcal{L}}
\newcommand{\A}{\mathcal{A}}
\newcommand{\K}{\mathcal{K}}
\newcommand{\B}{\mathcal{B}}
\newcommand{\D}{\mathcal{D}}
\newcommand{\ER}{\mathcal{R}}
\newcommand{\Ff}{\mathcal{F}}
\renewcommand{\P}{\mathcal{P}}
\newcommand{\G}{\mathcal{G}}
\newcommand{\Cc}{\mathcal{C}}
\newcommand{\E}{\mathcal{E}}
\newcommand{\NN}{\mathcal{N}}
\newcommand{\In}[1]{\text{\textcolor{blue}{\texttt{In[#1]}}}}
\DeclarePairedDelimiter\jap{\langle}{\rangle}
\newcommand{\fq}{F_{\alpha_4}}
\renewcommand{\phi}{\varphi}
\newcommand{\dx}{\partial_x}
\newcommand{\dxx}{\partial_{xx}}
\newcommand{\dt}{\partial_t}
\newcommand{\dtt}{\partial_{tt}}
\newcommand{\F}{\mathcal{F}}
\newcommand{\h}{\tilde{H}}
\newcommand{\ham}{\mathcal{H}}
\renewcommand{\H}{\mathcal{H}}
\newcommand{\supp}{{\rm supp}}
\newcommand{\bal}{{\bf \alpha}}
\newcommand{\bbt}{{\bf \beta}}
\newcommand{\I}{{\texttt{\bf J}}}
\newcommand{\J}{{\texttt{\bf J}}}
\newcommand{\Id}{{\rm Id}}
\newcommand{\fo}{F^1_{\alpha_4}}
\newcommand{\ft}{F^2_{\alpha_4}}
\newcommand{\ga}{G_{\alpha_2}}
\newcommand{\eb}{\bar{\eta}}
\newcommand{\p}{\psi}
\newcommand{\pb}{\bar{\psi}}
\newcommand{\w}{\bar{w}}
\newcommand{\W}{\bar{W}}
\newcommand{\z}{\bar{z}}
\newcommand{\BB}{B_\dc(\T,\C)}
\newcommand{\xt}{X_{{\bf Res},3}}
\newcommand{\xb}{\bar{\xi}}
\renewcommand{\epsilon}{\varepsilon}
\newcommand{\y}{\texttt{y}}
\newcommand{\Mat}{\texttt{M}}
\renewcommand{\Re}{{\rm Re}}
\renewcommand{\u}{\bar{u}}
\newcommand{\FO}{\mathcal{F}}
\newcommand{\fk}{\FO_{\bf kir}}
\newcommand{\mk}{\M_{\bf kirk}}
\newcommand{\ad}{{\rm ad}}
\newcommand{\id}{{\rm Id}}
\newcommand{\rg}{{\rm Rg}}
\renewcommand{\S}{\mathcal{S}}
\newcommand{\Rn}{\R^{n}}
\newcommand{\cross}[1]{\int_{\Rn}\nabla^{#1}u\cdot\nabla^{#1}u_t\,dx}
\newcommand{\nrmu}[1]{\lVert\nabla^{#1}u\rVert^{2}}
\newcommand{\nrmut}[1]{\lVert\nabla^{#1}u_t\rVert^{2}}
\renewcommand{\qedsymbol}{$\blacksquare$}
\renewcommand{\ge}{\geq}
\newcommand{\be}{{\bf e}}
\newcommand{\cF}{{\mathcal F}}
\newcommand{\im}{{\rm i}}
\newcommand{\cM}{{\mathcal M}}
\newcommand{\cR}{{\mathcal R}}
\newcommand{\cK}{{\mathcal K}}
\newcommand{\cG}{{\mathcal G}}
\newcommand{\mic}[1]{\textcolor{blue}{#1}}
\newcommand{\simo}[1]{\textcolor{violet}{#1}}
\newcommand{\td}{{\mathtt{d}}}
\newcommand{\tn}{{\mathtt{n}}}
\newcommand{\set}[1]{{\left\{#1\right\}}}
\newcommand{\bcoeffu}[1]{#1_{\bal^{1},\bbt^{1}}}
\newcommand{\bcoeffd}[1]{#1_{\bal^{2},\bbt^{2}}}
\newcommand{\buno}{b^{1}}
\newcommand{\bdue}{b^{2}}
\newcommand{\aluno}{\alpha^{1}}
\newcommand{\aldue}{\alpha^{2}}
\newcommand{\btuno}{\beta^{1}}
\newcommand{\btdue}{\beta^{2}}
\newcommand{\uu}{{u^\alpha}{\bar{u}^\beta}}
\newcommand{\uuu}{u^{\aluno}\bar{u}^{\btuno}}
\newcommand{\uud}{u^{\aldue}\bar{u}^{\btdue}}
\newcommand{\buu}{{u^\bal}{\bar{u}^\bbt}}
\newcommand{\buuu}{u^{\baluno}\bar{u}^{\bbtuno}}
\newcommand{\buud}{u^{\baldue}\bar{u}^{\bbtdue}}
\newcommand{\pon}{{{\Pi^{0,\mathcal{K}}}}} %%%%% prima era \cN ma ora con le nuove notazioni e la storia del Taylor ? proprio K
\newcommand{\por}{{\Pi^{0,\mathcal{R}}}}
\newcommand{\pd}{{\Pi^{-2}}}
\newcommand{\ono}[1]{#1^{0,\mathcal{N}}}
\newcommand{\oro}[1]{#1^{0,\mathcal{R}}}
\newcommand{\odo}[1]{#1^{-2}}
\newcommand{\modi}[1]{\norm{u_{#1}}^2}
\newcommand{\modu}{\norm{u}^2}
\newcommand{\cI}{{\mathcal I}}
\newcommand{\La}[1]{\Lambda^{\pa{#1}}}
%\newcommand{\es}{e^{\set{S,\cdot}}}
%\newcommand{\bal}{{ \al}}
%\newcommand{\bbt}{{ \bt}}
%\newcommand{\baluno}{\bal^{(1)}}
%\newcommand{\baldue}{\bal^{(2)}}
%\newcommand{\bbtuno}{\bbt^{(1)}}
%\newcommand{\bbtdue}{\bbt^{(2)}}
%\newcommand{\Heta}{\cH_{r,p,s,a,\eta}}
%\newcommand{\peseta}{e^{\eta\norm{\pi(\al - \bt)}}} %%%%% peso del momento 
%\newcommand{\nore}[1]{\norm{#1}_{r,p,s,a,\eta}}
%\newcommand{\noreq}[1]{\norm{#1}_{r,p,\eta}^{\wc} } %%% norma s,r,eta
%\newcommand{\Lnor}[1]{\norm{#1}^{\wc, \rm Lip}}  %%% seminorma lipschitz
%\newcommand{\gnor}[1]{\norm{#1}^{\wc,\gamma}}  %%% norma gamma
%\newcommand{\Dom}[1]{\sum_{#1} \omega_{#1} \pa{\norm{u_{#1}}^2 - I_{#1}}}  %% D_omega con l'indice che voglio
%\newcommand{\azioj}[1]{\pa{\norm{u_{#1}}^2 - I_{#1}}} %%% coef azioni con indici 
%\newcommand{\azio}{\pa{\norm{u}^2 - I}} %% coeff azione
%\newcommand{\divisor}[1]{#1\cdot\pa{\bal - \bbt}}
%\newcommand{\zero}[1]{#1^{(0)}}
%\newcommand{\zeroR}[1]{#1^{(0,\cR)}}
%\newcommand{\zeroK}[1]{#1^{(0,\cK)}}
%\newcommand{\due}[1]{#1^{(-2)}}
%\newcommand{\buon}[1]{#1^{\ge 2}}
%\newcommand{\costa}{C e^{C(\teta,p){\pa{\frac{\sigma}{2}}}^{-3/\teta}}}  %% la costante del lemma eqz homologica
%\newcommand{\call}{\mathcal{L}}
%\newcommand{\esp}[2]{e^{\set{#1, #2}}}


%\newcommand{\q}{\tilde{\mathcal{Q}}}
\newcommand{\q}{\mathcal{Q}}
%\newcommand{\q}{\dfrac{1}{2}\mu\left(Q(z,\z)\right)}
\newcommand{\dc}{\delta_0}
\renewcommand{\c}{\mathfrak{c}}
\newcommand{\rh}{{\bf \varrho}}
\newcommand{\qb}{\bar{q}}
\newcommand{\kk}{\K_{\bf kir}}
\newcommand{\rk}{\mathcal{R}_{\bf kir}}
\newcommand{\ZZ}{\mathcal{Z}}
\newcommand{\gF}{\mathscr{F}}
\renewcommand{\a}{\alpha}
\newcommand{\U}{\mathcal{U}}
\newcommand{\V}{\mathcal{V}}
\newcommand{\g}{\mathfrak{g}}
\newcommand{\con}{\mathfrak{C}}
\renewcommand{\O}{\mathcal{O}}
\newcommand{\Ra}{\mathcal{R}}
\newcommand{\Ka}{\mathcal{K}}
\newcommand{\QQ}{\mathcal{Q}}
\newcommand{\n}{\hat{n}}
\renewcommand{\l}{\ell}
\newcommand{\f}{{\bf f}}
\renewcommand{\d}{{\mathtt D}_\gamma}
\renewcommand{\k}{{\bf k}}
\renewcommand{\j}{{\bf j}}
\newcommand{\bh}{{\bm h}}
\newcommand{\zh}{ {\bf z}^{(\bh)}}
\newcommand{\al}{{\alpha}}
\newcommand{\bt}{{\beta}}
\newcommand{\pa}[1]{{\left(#1\right)}}
\newcommand{\bcoef}[1]{#1_{\bal,\bbt}}
\newcommand{\croc}{{\mathtt C}}
\newcommand{\s}{\sigma}
\newcommand{\kap}{\kappa}
\newcommand{\es}[1]{e^{\{\cdot,S_{#1}\}}}
\renewcommand{\gg}{\bf g}
\renewcommand{\epsilon}{\varepsilon}
\newcommand{\op}{\mathfrak{b}}
\newcommand{\OP}{\mathfrak{B}}
\newcommand{\TT}{\mathcal{T}}
\newcommand{\qs}{\q^{\bf Sym}}
\newcommand{\ds}{\mathtt{D}^{{\bf Sym}}_\gamma}
\newcommand{\m}{\hat{m}}
\newcommand{\tom}{\tilde{\omega}}
\newcommand{\Lambd}{\texttt{\bf M}}
\newcommand{\lambd}{\texttt{\bf m}}
\newcommand{\Hk}{\texttt{\bf H}}
\newcommand{\HK}{\mathscr{H}}
\newcommand{\Km}{\texttt{K}}
\newcommand{\dss}{\mathtt{D}_{\gamma,\S}}
\newcommand{\ww}{\texttt{w}}
\newcommand{\Ip}{\mathcal{I}}
\newcommand{\IM}{\bm{\mathcal{I}}}
\newcommand{\eps}{{\epsilon}_*}
\newcommand{\hame}{\hat{\ham}^{(\epsilon)}_{\scriptscriptstyle \bf BNF}}
\renewcommand{\le}{\leqsssim}
\renewcommand{\ge}{\gtrsim}
\newcommand{\ddt}{\dfrac{d}{dt}}
\providecommand{\vect}[2]{{\bigl(\begin{smallmatrix}#1\\#2\end{smallmatrix}\bigr)}} \providecommand{\sign}{\mathrm{sgn}\,}  
\providecommand{\sm}[4]{{\bigl(\begin{smallmatrix}#1&#2\\#3&#4\end{smallmatrix}\bigr)}}
\newcommand{\Ph}{\mathfrak{\Phi}}
\newcommand{\red}[1]{\textcolor{red}{#1}}
\newenvironment{multisys}
{\begin{minipage}{\displaywidth}
	\lineskip=1ex
	\leftskip=0pt plus 1fil
	\rightskip=.15\displaywidth plus 1fill
	\parfillskip=-.15\displaywidth plus -1fill
	$}
{$\end{minipage}}

\newenvironment{system}
{\left\lbrace\renewcommand{\arraystretch}{1.2}\begin{array}{@{}l@{}}}
{\end{array}\right.\kern-\nulldelimiterspace}
\newenvironment{claim}[1]{\par\noindent\underline{Claim:}\space#1}{}
\newenvironment{claimproof}[1]{\par\noindent\underline{Proof:}\space#1}{\hfill $\blacksquare$}
\newcounter{mathematicapage}
%\setlength{\parindent}{0cm}


\setcounter{tocdepth}{1}
\makeatletter

%%Add bold to section

%\renewcommand{\tocsection}[3]{%
%  \indentlabel{\@ifnotempty{#2}{\bfseries\ignorespaces#1 #2\quad}}\bfseries#3}
%


%%Allinea le sez e le sottosez
%
\def\l@subsection{\@tocline{2}{0pt}{2.5pc}{5pc}{}}
\def\l@subsubsection{\@tocline{3}{0pt}{4.5pc}{5pc}{}}
%Make Chapter disapear in ToC
\renewcommand\tocchapter[3]{%
\indentlabel{\@ifnotempty{#2}{\ignorespaces#2.\quad}}#3%
}
%\newcommand\@dotsep{4.5}
%\def\@tocline#1#2#3#4#5#6#7{\relax
%	\ifnum #1>\c@tocdepth % then omit
%	\else
%	\par \addpenalty\@secpenalty\addvspace{#2}%
%	\begingroup \hyphenpenalty\@M
%	\@ifempty{#4}{%
	%		\@tempdima\csname r@tocindent\number#1\endcsname\relax
	%	}{%
	%		\@tempdima#4\relax
	%	}%
%	\parindent\z@ \leftskip#3\relax \advance\leftskip\@tempdima\relax
%	\rightskip\@pnumwidth plus1em \parfillskip-\@pnumwidth
%	#5\leqavevmode\hskip-\@tempdima{#6}\nobreak
%	\leqaders\hbox{$\m@th\mkern \@dotsep mu\hbox{.}\mkern \@dotsep mu$}\hfill
%	\nobreak
%	\hbox to\@pnumwidth{\@tocpagenum{#7}}\par
%	\nobreak
%	\endgroup
%	\fi}
%\makeatother
%\AtBeginDocument{%
%	\makeatletter
%	\expandafter\renewcommand\csname r@tocindent0\endcsname{0pt}
%	\makeatother
%}
\def\l@subsection{\@tocline{2}{0pt}{2.5pc}{5pc}{}}
\newcommand\atotoc[1]{\addtocontents{toc}{#1\par}}
\title{On the integrability of the  Kirchhoff-Pohozaev equation on tori}
\author{Dario Bambusi}
\address{\scriptsize{Dipartimento di Matematica, Universit\`a degli Studi di  Milano, 
	Via Saldini 50: 20133 Milano }}
\email{dario.bambusi@unimi.it}
	\author{Emanuele Haus}
	\address{\scriptsize{Dipartimento di Matematica e Fisica, Universit\`a degli Studi RomaTre, 
	Largo San Leonardo Murialdo 1, 00146 Roma}}
	\email{emanuele.haus@uniroma3.it}
	\author{Simone Marrocco}
	\email{simone.marrocco@uniroma3.it}
	\author{Michela Procesi}
	\email{michela.procesi@uniroma3.it}
	
	\begin{document}
\maketitle
\begin{abstract}
In this paper we study the Kirchhoff-Pohozaev equation introduced in \cite{P2} and its constants of motion (see \cite{BoitiManfrin2025}, \cite{BoitiManfrin2026}) on $n$ dimensional tori. We show that these Hamiltonians are all in involution and we prove that they are generated by an infinite list of constants of motion which are all defined and in involution on a fixed phase space.  Then we study the Kirchhoff-Pohozaev equation restricted to a finite Fourier support. In dimension $n=1$ we show that such finite dimensional reduction is always completely integrable and provide an analytic Brikhoff normal form in a neighborhood of the origin. We also give sufficient conditions for integrability for $n>1$. We finally show that the formal Birkhoff Normal form of the Kirchhoff-Pohozaev equation is integrable for $n=1$.%, which we denote by $\cI^{(1)}_\nu,\cI^{(2)}_\nu$ where $\nu$ runs over the eigenvalues of the Laplacian, 
\end{abstract}
	\tableofcontents 
\section{Equation and main results}
The Kirchhoff-Pohozaev equation on the $n$-dimensional torus $\T^n$ is given by
\begin{equation}\label{KP-orig}
	u_{tt}-\left(\dfrac{1}{a+b\int_{\T^n}|\nabla u|^2dx}\right)^2\Delta u=0, \qquad (t,x) \in \R\times\T^n, 
\end{equation}
where $a$ and $b$ are nonzero real constants. By suitably scaling time and the size of the solution $u(t,x)$, the constants $a$ and $b$ can be normalized to $1$, apart from the sign of $b$. For the sake of simplicity, throughout our work we set $$a = b = 1,$$
namely we  consider the equation
\begin{equation}\label{KP}
	u_{tt}-\left(\dfrac{1}{1+\int_{\T^n}|\nabla u|^2dx}\right)^2\Delta u=0, \qquad (t,x) \in \R\times\T^n.
\end{equation}
As phase space we consider 
\begin{equation}\label{phase space}
	{\mathcal H}^s:= H^s(\T^n) \times H^{s-1}(\T^n)\,,\quad s\geq 1\,,
\end{equation}
endowed with the norm
\begin{equation}\label{normuv}
	\|(u,v)\|_{\H^s}^2:= \|u\|_{H^s}^2+  \|v\|_{H^{s-1}}^2\,.
\end{equation}
The Kirchhoff-Pohozaev equation \eqref{KP-orig} was first introduced by Pohozaev \cite{P2} as an example of a Kirchhoff-type equation with a higher order conservation law, which also allowed him to prove the global well-posedness of the equation.
Very recently, Boiti and Manfrin (see \cite{BoitiManfrin2025}, \cite{BoitiManfrin2026}) proved that equation \eqref{KP-orig} admits infinitely many conservation laws, at every order.
This, of course, raises some deep questions about the possible existence of an integrable structure of the equation, especially in space dimension $n=1$. In this context we prove both  formal integrability for the \eqref{KP} equation and analytic integrability for its finite dimensional reductions.
Let us give an informal overview.

As a first step we  give a quite explicit formulation of the constants of motion (see Proposition \ref{costanti!}) and then prove that, in any space dimension $n$, the infinitely many constants of motion found by Boiti and Manfrin are in involution, see Proposition \ref{involuz} . To do so we construct a list of constants of motion, which we denote $(\I_k)_{k\geq 2}$ which control the Sobolev norms and such that the constants of motion of Boiti and Manfrin are all polynomials in the $\J_k$. This procedure is quite abstract and works in more general settings, such as equation \eqref{KP} on compact manifolds without boundary. We mention the  preprint \cite{campos} (which was posted on arxiv  when this manuscript  was almost completed) for a similar result in a slightly different setting, see the end of the Introduction for more detailed remarks. 



The next step, which is our first main result  Theorem\ref{genero}, is to introduce  a new set of constants of motion, which we denote by $\cI_{\nu}^{(a)}$ with $a=1,2$ and $\nu$ running over the eigenvalues of the Laplacian. These constants of motion are all defined on the energy space $\H^1$ , and are in involution   in $\H^3$. Moreover we show that all the $\I_k$ as well as the Kirchhoff-Pohozaev Hamiltonian are expressed as absolutely convergent sums of the  $\cI_{\nu}^{(a)}$. Here we use the fact that the space domain is a torus $\T^n$.
In conclusion we have constructed a commuting family of Hamiltonians in $\H^3$, we call this the Kirchhoff-Pohozaev hierarchy.

Then we study the Kirchhoff-Pohozaev hierarchy resticted to a  finite Fourier support, which is well known to be an invariant subspace for the dynamics. We provide sufficient conditions on the support which guarantee that the restricted finite-dimensional system is integrable. In particular in our second main result Theorem \ref{finito}, we prove that in space dimension $n=1$  the Kirchhoff-Pohozaev hierarchy resticted to any  finite Fourier support is integrable and, in a neighborhood of zero  can be conjugated to a normal form depending only  on linear actions $p_j^2+q_j^2$.

We conclude by proving in Theorem \ref{formal-birk}, that  the Kirchhoff-Pohozaev hierarchy in dimension $n=1$  is formally integrable, namely there exists a formal symplectic change of variables which conjugates all the Hamiltonians in the Kirchhoff-Pohozaev hierarchy to an integrable formal Birkhoff Normal form depending only  on linear actions $p_j^2+q_j^2$.

\bigskip

In order to give a precise statement of our main results let us introduce some notation.
Let us define the Hamiltonian function $H : {\mathcal H}^s \to \R$ by
\begin{equation}\label{hamiltonian1int}
	H(u,v):=\frac{1}{2}\int_{\T^n}v^2dx-\frac{1}{2}\left(\dfrac{1}{1+\int_{\T^n}\left|\nabla u\right|^2dx}\right).
\end{equation}
Then equation \eqref{KP} is equivalent to the first-order system
\begin{equation}\label{sysJ}
	\dt\begin{pmatrix}
		u\\
		v
	\end{pmatrix}=J
	\begin{pmatrix}
		\nabla_u H\\
		\nabla_v H
	\end{pmatrix},
\end{equation}
where $\nabla_u, \nabla_v$ denote the $L^2$-gradients and $J$ is the Poisson tensor
\begin{equation}\label{J}
	J=\begin{pmatrix}
		0&I\\
		-I&0
	\end{pmatrix}.
\end{equation}
%This obviously implies that, if $u(t,x)$ solves \eqref{KP}, then the quantity
%\begin{equation}\label{hamiltonian ut}
%			H(u,u_t)=\frac{1}{2}\int_{\T^n}u_t^2dx-\frac{1}{2}\left(\dfrac{1}{1+\int_{\T^n}\left|\nabla u\right|^2dx}\right)
%\end{equation}
%is a constant of motion.
Given two functions $F,G : {\mathcal H}^s \to \R$, their Poisson bracket is defined by
\begin{equation}\label{def Poisson}
	\{F,G\}:= \int_{\T^n} \left( \nabla_u F \nabla_v G - \nabla_v F \nabla_u G \right)\,dx.
\end{equation}


Now we use  the fact that the space domain is the  torus $\T^n$ and  write $u$ and $v$ in Fourier series
\begin{equation}\label{fourier}
	u(t,x) = \sum_{j\in\Z^n\setminus\{0\}} u_j(t) \frac{e^{ij \cdot x}}{(2\pi)^\frac{n}{2}}, \qquad v(t,x) = \sum_{j\in\Z^n\setminus\{0\}} v_j(t) \frac{e^{ij \cdot x}}{(2\pi)^\frac{n}{2}}.
\end{equation}
We define 
\begin{equation}
	\Lambda = \left\{ \nu \in \mathbb{R}^+ : \nu = |j|^2 \text{ for some } j \in \mathbb{Z}^n \setminus \{0\} \right\}\,.
\end{equation}
For any $\nu \in \Lambda$, let us denote  $S_\nu:= \{j\in \Z^n\setminus \{0\} :  |j|^2=\nu\}$.
We define:
\begin{equation}\label{eq:Gamma1}
	\mathcal I_\nu^{(1)} = \sum_{j \in S_\nu} \left( A \nu |u_j|^2 + |v_j|^2 + \sum_{\Lambda\ni \nu'\ne \nu} \sum_{k \in S_{\nu'}} \frac{\nu \nu' }{\nu - \nu'} |u_j v_k - u_k v_j|^2 \right)
\end{equation}
\begin{equation}\label{eq:Gamma2}
	\mathcal I_\nu^{(2)} = \frac{1}{2} \sum_{j, k \in S_\nu} |u_j v_k - u_k v_j|^2\,,\qquad M_j:=  \frac{1}{2}  (u_j v_{-j} - u_{-j}v_j)\,.
\end{equation}

\begin{theorem}\label{genero}
	For all $\nu\in \Lambda$ the functions $\mathcal I^{(n)}_\nu$ with $n=1,2$  are bounded (and hence analytic) polynomials in $(u,v)\in \H^1$. Moreover for $(u,v)\in \H^3$ we have  
	\[
	\{\mathcal I^{(n)}_\nu,\mathcal I^{(m)}_{\nu'} \} = 	\{\mathcal I^{(n)}_\nu, M_j\}=0 \,,\qquad \forall n,m=1,2,\ \forall \nu,\nu'\in \Lambda, \ \forall j\in \mathbb{Z}^n \setminus \{0\}.
	\]
	Finally
	
	\begin{equation}\label{calIJ}
		H = \frac12(\sum_{\nu \in \Lambda} \mathcal I^{(1)}_\nu-1\,),\quad
		\I_k = \sum_{\nu \in \Lambda} \nu^{k-1} \mathcal I^{(1)}_\nu + (k-1) \sum_{\nu \in \Lambda} \nu^{k} \mathcal I^{(2)}_\nu.
	\end{equation}
	where the first sum is absolutely convergent in $\H^{2}$ while the second is so in  $\H^{k}$.
\end{theorem}

\begin{remark}
	Note that the  $\I_k$ are the constants of motion in formula \eqref{costanti!}
 and the formula \eqref{calIJ} may appear  quite surprising. In fact at a purely formal level, they are deduced in a very natural way from the Taylor expansion of the generating function of the $(\I_k)_{k\geq 2}$.
\end{remark}
%{\bf Integrability.} Recognizing and studying integrable systems is a fundamental problem in Hamiltonian dynamics both in finite and infinite dimension.
%\\
%For an analytic  Hamiltonian system on a $2k$ dimensional phase space,  the key point is to prove  the existence of $k$ analytic constants of motion $H_1,\dots,H_k$ which are generically independent and in involution. Furthermore one often can ensure that  the $H_i$ have a common elliptic fixed point (say at zero), with independent quadratic parts, Then one has:
%
%\begin{itemize}
%	\item  the   {\it Arnold Liouville} Theorem, which  guarantees that  a neighborhood of any regular and  compact  level set is foliated by Lagrangian tori,  parametrized by symplectic action-angle coordinates $(I_i,\theta_i)_{i=1}^k$, on  which the dynamics is conjugated to the linear  traslation $I(t)= I(0)$, $\theta(t)=  \theta(0) +\omega(I(0)) t$.
%	\item  the {\it Vey theorem} which ensures the existence of a convergent   Birkhoff Normal Form  change of variables, namely a symplectic analytic and invertible map defined in some ball at the origin,  so that in the new variables $(p_i,q_i)_{i=1}^k$ the linear actions $p_i^2+q_i^2$ are constants of motion(see \cite{Vey1978}, \cite{Eliasson1990},\cite{Ito1989}\cite{Stolovitch2000}, \cite{Zung2005}).
%\end{itemize} 
%In the infinite dimensional setting the picture is more complicated and there are a number of open problems. In the context of PDEs on tori, a main approach is to prove the existence of a {\it Lax pair} (see  \cite{Lax1968INTEGRALSON},\cite{Zakharov1971KortewegdeVE}) which automatically ensures the existence of an infinite family of constants of motion. The next step is to use this very strong structure to prove that:
%
%\begin{itemize}
%	\item  generic bounded solutions are almost periodic, in the sense that they are uniform limit of time quasi-periodic solutions.
%	\item  there exists a Birkhoff map, namely for all $s>s_0$, a symplectic analytic and invertible map defined in some ball at the origin in $\H^s$,  so that in the new variables $(p_i,q_i)$ the linear actions $p_i^2+q_i^2$ are constants of motion.
%\end{itemize} 
%This strategy has been carried out successfully for various PDEs on the circle such as the KdV, NLS, Benjamin-Ono equations.  in a more abstract setting an infinite dimensional version of Vey theorem was disucssed in \cite{KuksinPerelman2010} and \cite{BambusiStolovitch2020}.

A natural question is whether the  constants \eqref{eq:Gamma1}, \eqref{eq:Gamma2} are related to an integrable structure at least locally close to zero, following the approach of \cite{KuksinPerelman2010} or \cite{BambusiStolovitch2020}. Note that in both these results one needs a set of constants of motion all defined on the same phase space. This was in fact a motivation for Theorem \ref{genero} however we remark that our constants of motion do not fit their hypotheses. To partially bypass this problem we use the 
 well known fact  that  Kirchhoff type equations preserve the Fourier support.
\\
Given any symmetric\footnote{Namely such that if $j\in T$ then $-j\in T$. We need this condition in order to find real solutions.} set $T\subset \Z^n\setminus\{0\}$ we have that the subspace
\begin{equation}
	\label{UT}
	\mathscr{U}_T:=\{(u,v)\in \H^1 \;\vert{} \quad u_j=v_j=0 \,,\quad \forall j\notin T \}
\end{equation}
is invariant under the dynamics.
As a consequence, if $T$ is finite, we obtain a finite dimensional Hamiltonian system with $|T|$ degrees of freedom, so it is very natural to wonder whether and how many of  the constants of motion $\cI,M$  survive this reduction.  In our setting the key point is the decomposition of $T$ in eigenspaces of the Laplacian. 	We may write, w.l.o.g, $T=\cup_{h=1}^N T_h$ where $T_h= T\cap S_{\nu_h}\neq \emptyset$ for some  increasing sequence $\nu_h\in \Lambda$. We prove (see Lemma \ref{restringo}) that the number of algebraically independent constants of motion $\cI,M$ defined in \eqref{eq:Gamma1}, \eqref{eq:Gamma2} is at least $N+|T|/2$ and at most
$
N + K+ |T|/2   
$
where $K$ is the number of $T_h$ with cardinality $\geq 4$.
\\
In particular in dimension $n=1$ the number of independent constants of motion is equal to $|T|$ and the restricted system is integrable. 
More precisely we have the following
\begin{theorem}\label{finito}
	Consider any symmetric finite set $T=\{j_1,-j_1,\dots, j_d,-j_d\}\subset  \Z^n\setminus\{0\}$, with $|j_l| \neq | j_m|$ if $l\neq m$. Then $\mathscr{U}_T\equiv  \R^{4d}$ and:
	
	$(i)$ \ the constants of motion
	\begin{equation}
		\label{massimali?}
		\cI^{(1)}_{|j_1|^2},\dots, \cI^{(1)}_{|j_d|^2}\,,\; M_{j_1},\dots M_{j_d}
	\end{equation}
	are a set of $2d$ algebraically independent constant of motion in involution, which have an elliptic fixed point at $(u,v)=0$ with independent quadratic parts.
	
	
	$(ii)$ \ there exists $r=r(T)$ and a close to identity  symplectic analytic invertible map $\Phi: B_r(\R^{4d}) \to \R^{4d}$, $(u,v)\mapsto (p,q)$ which preserves the $M_j$ and brings the $\cI^{(1)}_{|j_k|^2}$ (and thus $H$, via formula eqref{} ) in normal form $$\cI^{(1)}_{|j_k|^2}\circ \Phi=\mathscr{I}_{k}(\{p_{h}^2+q_{h}^2\}_{h=1,\dots,2d })\,.$$
\end{theorem}
Since for $n=1$ the condition  $|j_l| \neq | j_m|$ if $l\neq m$ is automatically met, any finite symmetric  set $T\in \Z\setminus\{0\}$ automatically satisfies the hypotheses of Proposition \ref{finito}. We thus have
\begin{corollary}\label{uffa}
	The Kirchhoff-Pohozaev equation on the circle restricted to any finite dimensional set $\mathscr{U}_T$ is integrable. Moreover there exists a neighborhood of zero where the Birkhoff Normal Form algorithm converges and solutions all live on tori.
\end{corollary}
\begin{remark}
	We point out that the condition $|j_i|\neq |j_m|$ in Theorem \ref{finito} is quite important. To see this consider  for $d\geq 2$,
	$T= \{j_1,-j_1,\dots, j_d,-j_d\}\subset  \Z^n\setminus\{0\}$, with  the $j_i$ distinct but with $|j_l| = | j_m|$ for all $l,m=1,\dots,d$. Of course this may occur only if the space dimension $n>1$.  Direct inspection of Formulas \eqref{eq:Gamma1}, \eqref{eq:Gamma2} shows that at most $d+2$ constants are non identically zero. If $d>2$ then $d+2<2d$.
\end{remark}

\smallskip

Of course the integrability of the restricted system  for $n=1$ does not imply that one is able to integrate the Kirchhoff-Pohozaev equation even close to zero.
The information which we have from the constants of motion in involution is however sufficient to prove, in the case $n=1$, the integrability of the formal Birkhoff normal form. More precisely we have the following.

\begin{theorem}\label{formal-birk}
	There exists a formal symplectic change of variables which preserves the $M_j$ and conjugates all the commuting Hamiltonians $\cI_k$, with $k\geq 1$ (and consequently all the $\I_k$ and $H$ )  to a normal form  $\mathcal N_k(\{p_j^2+q_j^2\}_{j\in \Z\setminus\{0\}})$.
\end{theorem}
We conclude this introduction with a brief comparison with \cite{campos} which appeared on arxiv as we were completing our manuscript. The starting point of the two papers is  similar,  showing that there exists a sequence of constant of motion which at first order coincide with those of Boiti and Manfrin and are in involution. Then \cite{campos}  studies  a slightly more general model in which the Laplacian is substituted by a positive definite selfadjoint operator and focuses  on smoothness properties of the solutions proving equicontinuity in $H^s$. Here we focus on  integrability properties of the equation showing in particular that any finite dimensional reduction is integrable in the standard sense, at least in dimension 1. In this sense our approach allows to conclude that in a reasonable sense the Kirchoff Pohozaev equatin in dimension 1 is integrable, while the higher dimensional case is still open. 

\smallskip
\thanks{{\bf Acknowledgements.} We wish to thank Joackim Bernier and Claudio Procesi for helpul suggestions and discussions. DB was partially supported by GNFM, EH, SM and MP were partially supported by GNAMPA. 

\section{Constants of motion}

In this section, we show that equation \eqref{KP} admits infinitely many conserved quantities $\{\I_k\}_{k\geq2}$  of the form $\I_k=E_k+f(\{E_j, F_j, Q_j\})$, where $f$ is a homogeneous polynomial of degree $2$ in the variables $\{E_j, F_j, Q_j\}$ defined below in \eqref{BDQ}.

Following Boiti and Manfrin, for $k\in\N$ we start by setting the notation
$$
\nabla^{2k} u := \Delta^k u, \qquad \nabla^{2k+1} u := \nabla(\Delta^k u), \qquad \| u \| := \| u \|_{L^2(\T^n)} = \left( \int_{\T^n} |u|^2\, dx \right)^{\frac12}.
$$
Then, for $(u,v)\in{\mathcal H}^s$ and $k\in\N_{\geq1}$, we introduce the quantities (well-defined if $s\geq k$)
\begin{equation}\label{BDQ}
	\begin{aligned}
		D_k&:=\|\nabla^{k}u\|^2\quad\quad&
		B_k&:=\|\nabla^{k-1}v\|^2\quad\quad&
		Q_k&:=\int_{\T^n}\nabla^{k-1}u\cdot\nabla^{k-1}v\,dx\\
		A&:=\dfrac{1}{1+D_1}\quad\quad&
		E_k&:=\dfrac{B_k}{A}+A D_k\quad\quad&
		F_k&:=\dfrac{B_k}{A}-A D_k.
	\end{aligned}
\end{equation}
In this set of variables, the Hamiltonian takes the form
\begin{equation}\label{ham new}
	H=\frac12(B_1-A)=\frac12 AE_1 - A + \frac12 A^2.
\end{equation}
Also note that equation \eqref{KP} reads
\begin{equation}\label{KP A}
	u_{tt} = A^2(t) \Delta u,
\end{equation}
where $A=A(t)$ is the quantity defined in \eqref{BDQ}.
\begin{proposition}\label{costanti!}
	For $k\geq2$, the functions
	\begin{equation}\label{def I}
		\I_k= E_k- \sum_{j=2}^k Q_j Q_{k+2-j} + \frac14 \sum_{j=2}^{k-1}(E_j E_{k+1-j}-F_j F_{k+1-j})
	\end{equation}
	are independent constants of motion which are well defined on $\mathcal H^k$.
\end{proposition}

\begin{remark}
	The functions  $\I_k$ coincide with the $k$-th order conservation law found by Boiti and Manfrin at the highest order of derivation, but differ by lower order terms (which can be expressed as polynomials in the lower order constants of motion $\I_j$).  It is not hard to check that the constant $\I_k$  controls the Sobolev norm $\H^k$. In fact the $\I_k$ coincide with the constants of motion of \cite{campos}, see formula (3.1) of that paper. 
\end{remark}

\smallskip
In order to prove Proposition \ref{costanti!} we start by analyzing the time evolution of the building blocks \eqref{BDQ}.
\\
Given a function $F : {\mathcal H}^s \to \R$, we denote $\dot F = \{ F,H \}$, where $H$ is the Hamiltonian defined in \eqref{hamiltonian1int}.
In the next lemma, we compute the Poisson brackets of the quantities defined in \eqref{BDQ} with the Hamiltonian $H$.

\begin{lemma}\label{time derivatives}
	Let $D_k, B_k, Q_k, A, E_k, F_k$ be the quantities defined in \eqref{BDQ}. Then
	\begin{equation}\label{BDQderivate}
		\begin{aligned}
			\dot{D}_k&=2Q_{k+1}\quad\quad&
			\dot{B}_k&=-2A^2Q_{k+1}\quad\quad&
			\dot{Q}_k&=B_k-A^2D_k=AF_k\\
			\dot{A}&=-2A^2Q_2\quad\quad&
			\dot{E}_k&=2 AQ_2 F_k= 2Q_2\dot{Q}_k\quad\quad&
			\dot{F}_k&=2A\left(Q_2E_k-2Q_{k+1}\right).
		\end{aligned}
	\end{equation}
\end{lemma}
\begin{proof}
	Let $(u,v)$ solve \eqref{sysJ}. Then the formula
	\begin{equation}\label{dot D}
		\dot{D}_k=2Q_{k+1} 
	\end{equation}
	is obvious, since $u_t=v$. The formula for $\dot{B}_k$ follows immediately from \eqref{KP A} and integration by parts, since
	\begin{equation}\label{dot B}
		\begin{aligned}
			\dot{B}_k & = 2\int_{\T^n} \nabla^{k-1}v \cdot \nabla^{k-1}v_t\, dx = 2\int_{\T^n} \nabla^{k-1}v \cdot \nabla^{k-1}u_{tt}\, dx \\
			& = 2A^2\int_{\T^n} \nabla^{k-1}v \cdot \nabla^{k+1}u\, dx = - 2A^2\int_{\T^n} \nabla^{k}v \cdot \nabla^{k}u\, dx = -2A^2Q_{k+1}.
		\end{aligned}
	\end{equation}
	The formula for $\dot{Q}_k$ follows similarly, since
	\begin{equation}\label{dot Q}
		\begin{aligned}
			\dot{Q}_k & = \int_{\T^n} \nabla^{k-1}u_t \cdot \nabla^{k-1}v\, dx + \int_{\T^n} \nabla^{k-1}u \cdot \nabla^{k-1}v_t\, dx \\
			& = \| \nabla^{k-1}v \|^2 + \int_{\T^n} \nabla^{k-1}u \cdot \nabla^{k-1}u_{tt}\, dx
			= \| \nabla^{k-1}v \|^2 - A^2  \| \nabla^k u \|^2 = B_k-A^2D_k = AF_k.
		\end{aligned}
	\end{equation}
	Next, by \eqref{dot D} we have
	\begin{equation}\label{dot A}
		\dot A = -\frac{\dot{D}_1}{(1+D_1)^2} = -2A^2Q_2.
	\end{equation}
	Finally, using \eqref{dot D}-\eqref{dot A}, we get
	\begin{equation}\label{dot E}
		\dot{E}_k=\frac{\dot{B}_k}{A}-\frac{\dot A}{A^2}B_k+{\dot A}D_k+A\dot{D}_k= 2Q_2(B_k-A^2D_k)= 2AQ_2 F_k =2Q_2\dot{Q}_k
	\end{equation}
	and
	\begin{equation}\label{dot F}
		\dot{F}_k =\frac{\dot{B}_k}{A}-\frac{\dot A}{A^2}B_k-{\dot A}D_k-A\dot{D}_k
		= -4AQ_{k+1}+2Q_2(B_k+A^2 D_k) = 2A(Q_2E_k-2Q_{k+1}).
	\end{equation}
\end{proof}



\smallskip 
In order to prove Proposition \ref{costanti!} we prove two preliminary results.
\begin{lemma}
	We have that for all $2\leq  k_1\leq  k_2$
	\begin{equation}\label{mic}
		\begin{aligned}
			\dot{Q}_{k_1}Q_{k_2}=\begin{cases}
				\dfrac12 \ddt Q_{k_1}^2 \,,&\quad k_1= k_2 \\ &\\
				\dfrac18 \ddt (E_{k_1}^2 - F_{k_1}^2) \,, & \quad k_1= k_2-1\\ &\\
				\dfrac{1}{4}\ddt\left(E_{k_1}E_{k_2-1}-F_{k_1}F_{k_2-1}\right)-  Q_{k_1+1}\dot{Q}_{k_2-1}, &\quad 2\leq   k_1< k_2-1 
			\end{cases}
		\end{aligned}
	\end{equation}
\end{lemma}
\begin{proof}
	If $k_1=k_2$  \eqref{mic} is obvious. Otherwise, we solve for $Q_{k+1}$ in the formula for $\dot F_k$ in \eqref{BDQderivate} (applied with $k=k_2-1$) and we get
	\begin{equation}\label{mic2}
		\begin{aligned}
			\dot{Q}_{k_1}Q_{k_2}&= \dot Q_{k_1} \left(\dfrac{1}{2}Q_{2}E_{k_2-1}-\dfrac{\dot{F}_{k_2-1}}{4A}\right)= \frac14 \dot E_{k_1}E_{k_2-1} -\frac14 F_{k_1}\dot F_{k_2-1}
		\end{aligned}
	\end{equation}
	if $k_1= k_2-1$ this concludes the proof of \eqref{mic}, otherwise
	\[
	\begin{aligned}
		&\dot{Q}_{k_1}Q_{k_2}= \frac14 \dot E_{k_1}E_{k_2-1} -\frac14 F_{k_1}\dot F_{k_2-1} = \frac14\ddt ( E_{k_1}E_{k_2-1}-F_{k_1}F_{k_2-1}) -\frac14 E_{k_1}  \dot E_{k_2-1} +\frac14 \dot F_{k_1} F_{k_2-1}
		\\
		&= \frac14\ddt ( E_{k_1}E_{k_2-1}-F_{k_1}F_{k_2-1}) -\frac12 E_{k_1} Q_2 \dot Q_{k_2-1} +\frac 12 A( Q_2 E_{k_1} -2 Q_{k_1+1} ) F_{k_2-1}
		\\
		&= \frac14\ddt ( E_{k_1}E_{k_2-1}-F_{k_1}F_{k_2-1}) - Q_{k_1+1} \dot Q_{k_2-1}.
	\end{aligned}
	\]
\end{proof}

\begin{lemma}\label{provo}
	We have that for all $2\leq  k_1 \leq  k_2$
	\begin{equation}
		\dot{Q}_{k_1}Q_{k_2}=\ddt \left( \frac18 \sum_{j=0}^{k_2-k_1-1}(E_{k_1+j}E_{k_2-j-1}-F_{k_1+j}F_{k_2-j-1}) - \frac12 \sum_{j=1}^{k_2-k_1-1} Q_{k_1+j} Q_{k_2-j} \right).
	\end{equation}
\end{lemma}
\begin{proof}
	We proceed by induction on $k_2-k_1$. If $k_2-k_1=1$, the thesis follows immediately by the second line in \eqref{mic}.
	If $k_2-k_1\geq2$ we use the third line in \eqref{mic} and obtain
	\begin{equation}\label{Q dot Q}
		\begin{aligned}
			\dot{Q}_{k_1}Q_{k_2}
			&= \frac14\ddt ( E_{k_1}E_{k_2-1}-F_{k_1}F_{k_2-1}) - Q_{k_1+1} \dot Q_{k_2-1}
			\\
			&= \ddt\left( \frac14( E_{k_1}E_{k_2-1}-F_{k_1}F_{k_2-1}) - Q_{k_1+1} Q_{k_2-1}\right) + \dot Q_{k_1+1} Q_{k_2-1}
		\end{aligned}
	\end{equation}
	Since $k_1+1\leq k_2-1$, we can apply the induction hypothesis to $\dot Q_{k_1+1} Q_{k_2-1}$. Thus, we get from \eqref{Q dot Q}
	\begin{equation}\label{Q dot Q 2}
		\begin{aligned}
			&\dot{Q}_{k_1}Q_{k_2}
			= \ddt\left( \frac14( E_{k_1}E_{k_2-1}-F_{k_1}F_{k_2-1}) - Q_{k_1+1} Q_{k_2-1}\right)\\
			& + \ddt \left( \frac18 \sum_{j=0}^{k_2-k_1-3}(E_{k_1+j+1}E_{k_2-j-2}-F_{k_1+j+1}F_{k_2-j-2}) - \frac12 \sum_{j=1}^{k_2-k_1-3} Q_{k_1+j+1} Q_{k_2-j-1} \right)\\
			& = \ddt \left( \frac18 \sum_{j=0}^{k_2-k_1-1}(E_{k_1+j}E_{k_2-j-1}-F_{k_1+j}F_{k_2-j-1}) - \frac12 \sum_{j=1}^{k_2-k_1-1} Q_{k_1+j} Q_{k_2-j} \right).
		\end{aligned}
	\end{equation}
\end{proof}







\begin{corollary}\label{lemma1}
	Let $k\geq2$. We have that
	\begin{equation}
		\dot E_k =	 \ddt\left( \sum_{j=2}^k Q_2 Q_{k-j+2} - \frac14 \sum_{j=2}^{k-1}(E_j E_{k-j+1} - F_j F_{k-j+1}) \right).
	\end{equation}
\end{corollary}
\begin{proof}
	By \eqref{BDQderivate} we have $\dot E_k = 2Q_2 \dot Q_k$. If $k=2$, then $\dot E_2 = \ddt Q_2^2$, which is the thesis for $k=2$. If $k\geq3$, then
	\[
	\dot E_k = 2Q_2 \dot Q_k = 2\ddt(Q_2 Q_k)- 2 Q_2 \dot Q_k
	\]
	and the thesis follows by applying Lemma \ref{provo} with $k_1=2$, $k_2=k$.
\end{proof}
\begin{proof}[Proof of Proposition \ref{costanti!}]
	The fact that the $\I_k$ defined in \eqref{def I} are constants of motion is a direct consequence of Corollary \ref{lemma1}. The fact that they are independent  is due to the fact that the quadratic parts in $u,v$ are so.
\end{proof}


\begin{remark}
	To give some examples, we list explicitly the first few $\I_k's$:
	\begin{align*}
		\I_2 &= E_2 - Q_2^2 \\
		\I_3 &= E_3 - 2Q_2 Q_3 + \frac{1}{4}\left(E_2^2 - F_2^2\right) \\
		\I_4 &= E_4 - 2Q_2 Q_4 + \frac{1}{2}(E_2 E_3 - F_2 F_3) - Q_3^2 \\
		\I_5 &= E_5 - 2Q_2 Q_5 + \frac{1}{2}(E_2 E_4 - F_2 F_4) - 2Q_3 Q_4 + \frac{1}{4}\left(E_3^2 - F_3^2\right) \\
		\I_6 &= E_6 - 2Q_2 Q_6 + \frac{1}{2}(E_2 E_5 - F_2 F_5) + \frac{1}{2}(E_3 E_4 - F_3 F_4) - 2Q_3 Q_5 - Q_4^2
	\end{align*}
	%	\[
	%	\I_2 =E_2 - Q_2^2\,,\quad \I_3= E_3 -2 Q_2 Q_3 + \frac14(E_2^2 - F_2^2)
	%	\]
	%	\[
	%	\I_4=E_4-2Q_2Q_4 + \frac12(E_2 E_3 - F_2 F_3) - Q_3^2
	%	\]
	For instance, as one can explicitly verify, $\I_2$ coincides with the second order conservation law introduced by Pohozaev, while the third order conservation law found by Boiti and Manfrin is $\I_3-\frac14\I_2^2$ in our language.
\end{remark}
It may be convenient to define 
\begin{equation}\label{def:J1}
	\I_{1} = A - B_{1} - 1 = -2H -1.
\end{equation}


\section{Poisson algebra and involution}
The next step is naturally to investigate whether the constants of motion $\I_k$ are in involution. 
\begin{definition}
	Let $\mathscr{P}$ be the ring of finite\footnote{Since we have infinitely many variables we need to specify that we are here considering only finite sums of monomials.}  polynomials in the variables $\{E_k,F_k,Q_k\}_{k\geq 2}, A ,1/A$ defined in \eqref{BDQ}, namely of finite linear combinations of monomials
	\[
	P = \sum_{i=1}^N P_i A^{d_i}\,\prod_{k=2}^M E_k^{a_{i,k}} F_k^{b_{i,k}} Q_k^{c_{i,k}} \,,\quad a_{i,k},b_{i,k},c_{i,k}\in \N,\,d_i\in \Z\,,\quad  P_i\in \R\,.
	\]
	Let $\mathscr C$ be the subring of polynomials in $\mathscr P$ which are constants of motion for \eqref{KP}
	.\end{definition}

%\red{Moreover they all commute with the Kirchhoff constants of motion
	%\[
	%K_j= I_j-I_{-j}
	%\]}
%\red{come si scrivono nelle variabili $u,v$ forse da dire dopo?}
%
%\red{nelle variabili $u,v$ si ha $K_j= I_j-I_{-j}=\frac{i}{2}(v_j u_{-j}-u_j v_{-j})$. Ogni $K_j$ \`e conservata. Complessivamente, il momento \`e la quantit\`a (vettoriale se siamo in dimensione $n\geq2$)
	%$$
	%M = \int_{\T^n} u\nabla v\, dx = \sum_{j\in\Z^n\setminus\{0\}} j K_j = \frac{i}{2} \sum_{j\in\Z^n\setminus\{0\}} j (v_j u_{-j}-u_j v_{-j}).
	%$$
	%Chiaramente si ha $K_{-j}=-K_j$, quindi $jK_j = -jK_{-j}$, perci\`o la sommatoria \`e fatta da coppie di addendi uguali tra loro.
	%(rispetto al lavoro Baldi-Haus ho integrato per parti e definito il momento con il segno opposto, per evitare di cambiare il segno di $K_j$ o di avere uno sgradevole segno meno nella formula)
	%}
\begin{proposition}\label{involuz}
	The space $\mathscr{P}$ is a Poisson algebra, moreover $\mathscr{C}$ is generated\footnote{Namely any element of $\mathscr{C}$ is expressed as a polynomial in the $\J_k$.}  by  $(\I_k)_{k\geq 2}$.
	Finally, for any $k_1,k_2\geq 1$ we have 
	\[
	\{\I_{k_1},\I_{k_2}\} =0\,.
	\]
\end{proposition}
\begin{remark}
	We mention that the involution property of the $\I_k$ is proved also in \cite{campos}, see Proposition 3.2. Their proof relies on the construction of a generating function and on direct algebraic computations. Ours is more based on the Poisson algebra structure of $\mathscr{P}$. In fact it seems to us that our statement is slightly stronger, since Proposition \ref{involuz} amounts to saying that any constant of motion which is a finite polynomial in $\{E_k,F_k,Q_k\}_{k\geq 2}, A ,1/A$ can be expressed as a polynomial in the $\I_k$. This applies for instance to the  constants of motion of Boiti Manfrin.
\end{remark}


%\begin{definition}
%	Let $\mathscr{P}$ be the ring of finite  polynomials in the variables $\{E_k,F_k,Q_k\}_{k\geq 2}, A ,1/A$ defined in \eqref{BDQ}, namely of finite linear combinations of monomials
%	\[
%	P = \sum_{i=1}^N P_i A^{d_i}\,\prod_{k=2}^M E_k^{a_{i,k}} F_k^{b_{i,k}} Q_k^{c_{i,k}} \,,\quad a_{i,k},b_{i,k},c_{i,k}\in \N,\,d_i\in \Z\,,\quad  P_i\in \R\,.
%	\]
%	Let $\mathscr C$ be the subring of polynomials in $\mathscr P$ which are constants of motion for \eqref{KP}
%	.\end{definition}
%
%%\red{Moreover they all commute with the Kirchhoff constants of motion
%	%\[
%	%K_j= I_j-I_{-j}
%	%\]}
%%\red{come si scrivono nelle variabili $u,v$ forse da dire dopo?}
%%
%%\red{nelle variabili $u,v$ si ha $K_j= I_j-I_{-j}=\frac{i}{2}(v_j u_{-j}-u_j v_{-j})$. Ogni $K_j$ \`e conservata. Complessivamente, il momento \`e la quantit\`a (vettoriale se siamo in dimensione $n\geq2$)
%	%$$
%	%M = \int_{\T^n} u\nabla v\, dx = \sum_{j\in\Z^n\setminus\{0\}} j K_j = \frac{i}{2} \sum_{j\in\Z^n\setminus\{0\}} j (v_j u_{-j}-u_j v_{-j}).
%	%$$
%	%Chiaramente si ha $K_{-j}=-K_j$, quindi $jK_j = -jK_{-j}$, perci\`o la sommatoria \`e fatta da coppie di addendi uguali tra loro.
%	%(rispetto al lavoro Baldi-Haus ho integrato per parti e definito il momento con il segno opposto, per evitare di cambiare il segno di $K_j$ o di avere uno sgradevole segno meno nella formula)
%	%}
%\begin{proposition}\label{involuz}
%	The space $\mathscr{P}$ is a Poisson algebra, moreover $\mathscr{C}$ is generated by  $(\I_k)_{k\geq 2}$.
%	Finally, for any $k_1,k_2\geq 1$ we have 
%	\[
%	\{\I_{k_1},\I_{k_2}\} =0
%	\]
%\end{proposition}


We now compute the Poisson brackets between the building blocks defined in \ref{BDQ}.
\begin{lemma}\label{relations1}
The following  commutation relations hold:
\begin{equation}
	\begin{aligned}
		&\{D_{k_1},D_{k_2}\}=0,\quad\quad&\{B_{k_1},B_{k_2}\}=0\quad\quad&\{Q_{k_1},Q_{k_2}\}=0\\
		&\{D_{k_1},B_{k_2}\}=4Q_{k_1+k_2},\quad\quad&\{D_{k_1},Q_{k_2}\}=2D_{k_1+k_2-1}\quad\quad&\{B_{k_1},Q_{k_2}\}=-2B_{k_1+k_2-1}\\
		&\{A,D_{k}\}=0,\quad\quad&\{A,B_{k}\}=-4A^2Q_{k+1}\quad\quad&\{A,Q_{k}\}=-2A^2D_k
	\end{aligned}
\end{equation}
\end{lemma}
\begin{proof}
This is a simple direct computation.
\end{proof}
As a consequence we show that $\mathscr{P}$ is closed under Poisson brackets.
\begin{lemma}[Algebra Rules]\label{algebra}	We have the following commutation rules: 
\begin{align*}
	\{E_{k_1},E_{k_2}\} &= 4F_{k_1}Q_{k_2+1}-4F_{k_2}Q_{k_1+1} \,,\quad\qquad\quad\quad\;\;
	\{E_{k_1},Q_{k_2}\} = F_{k_1}\big(E_{k_2}-F_{k_2}\big)-2F_{k_1+k_2-1} 
	\\[4pt]
	\{E_{k_1},F_{k_2}\} &= 8Q_{k_1+k_2}+4F_{k_1}Q_{k_2+1}-4E_{k_2}Q_{k_1+1} \,,\quad 
	\{F_{k_1},F_{k_2}\} = 4E_{k_1}Q_{k_2+1}-4E_{k_2}Q_{k_1+1} 
	\\[4pt]
	\{F_{k_1},Q_{k_2}\} &= E_{k_1}\big(E_{k_2}-F_{k_2}\big)-2E_{k_1+k_2-1} 
\end{align*}
moreover
\begin{align}
	\{A,E_{k}\} &= \{A,F_{k}\} = -4A\,Q_{k+1} ,\quad 
	\{A,Q_{k}\} =  -A\big(E_{k}-F_{k}\big)\,. %\qquad \{E_k,F_k\} = 8\,Q_{2k}\,.
\end{align}
\end{lemma}
\begin{proof}
This is again a direct computation. Let us briefly prove the first equality. We expand by bilinearity:
\begin{equation}
	\{E_{k_1}, F_{k_2}\} = \left\{ \frac{B_{k_1}}{A}, \frac{B_{k_2}}{A} \right\} - \left\{ \frac{B_{k_1}}{A}, A D_{k_2} \right\} + \left\{ A D_{k_1}, \frac{B_{k_2}}{A} \right\} - \{A D_{k_1}, A D_{k_2}\}
\end{equation}

We evaluate the individual blocks using the Leibniz rule and the fundamental relations:
\begin{enumerate}
	%			\item \textbf{The $\{D, D\}$ block:} Since $\{A, D_k\} = 0$ and $\{D_{k_1}, D_{k_2}\} = 0$, it vanishes identically:
	%			\begin{equation}
		%				\{A D_{k_1}, A D_{k_2}\} = 0
		%			\end{equation}
	
	\item \textbf{The $\{B, B\}$ block:} Since $\{B_{k_1}, B_{k_2}\} = 0$ and $\left\{\frac{1}{A}, B_k\right\} = 4Q_{k+1}$, we obtain:
	\begin{equation}
		\left\{ \frac{B_{k_1}}{A}, \frac{B_{k_2}}{A} \right\} = \frac{B_{k_1}}{A}(4Q_{k_2+1}) - \frac{B_{k_2}}{A}(4Q_{k_1+1})
	\end{equation}
	
	\item \textbf{The cross $\{B, D\}$ blocks:} Computing the first term, given that $\{B_{k_1}, D_{k_2}\} = -4Q_{k_1+k_2}$ and $\left\{\frac{B_{k_1}}{A}, A\right\} = 4AQ_{k_1+1}$:
	\begin{equation}
		\left\{ \frac{B_{k_1}}{A}, A D_{k_2} \right\} = A\left\{\frac{B_{k_1}}{A}, D_{k_2}\right\} + D_{k_2}\left\{\frac{B_{k_1}}{A}, A\right\} = -4Q_{k_1+k_2} + 4AD_{k_2}Q_{k_1+1}
	\end{equation}
	By antisymmetry, swapping indices yields the second term:
	\begin{equation}
		\left\{ A D_{k_1}, \frac{B_{k_2}}{A} \right\} = 4Q_{k_1+k_2} - 4AD_{k_1}Q_{k_2+1}
	\end{equation}
	\item \textbf{The $\{D, D\}$ block:}  is clearly equal to zero.
\end{enumerate}

Substituting these four relations back into the initial expansion:
\begin{equation}
	\{E_{k_1}, F_{k_2}\} = 8Q_{k_1+k_2} + 4\left(\frac{B_{k_1}}{A} - AD_{k_1}\right)Q_{k_2+1} - 4\left(\frac{B_{k_2}}{A} + AD_{k_2}\right)Q_{k_1+1}
\end{equation}

Finally, recognizing the original definitions $F_{k_1} = \frac{B_{k_1}}{A} - AD_{k_1}$ and $E_{k_2} = \frac{B_{k_2}}{A} + AD_{k_2}$, we arrive at the general identity:
\begin{equation}
	\{E_{k_1}, F_{k_2}\} = 8Q_{k_1+k_2} + 4F_{k_1}Q_{k_2+1} - 4E_{k_2}Q_{k_1+1}
\end{equation}
\end{proof}
\begin{proof}[Proof of Proposition \ref{involuz}] The fact that $\mathscr P$ is closed under Poisson brackets is a direct consequence of Lemma \ref{algebra}. 
Now we note that $\I_k\in \mathscr C$ and moreover we may recusively invert each relation so that  that, for any $i\geq 2$
\[
E_i= P_i(\I_2,\dots,\I_i,F_2,\dots, F_{i-1},Q_2,\dots,Q_i)\,.
\]
for some polynomial $P_i$. 
Consider  a polynomial constant of motion
\[
g(Q_2,\dots, Q_{h}, F_2,\dots, F_{h}, E_2,\dots, E_h,,A,1/A)\in \mathscr C\,.
\]
Substituting the $E_i$, there must exist a finite polynomial  $f$ in the variables  $F,Q,\I,A,1/A$ such that 
\[
\begin{aligned}
	f(Q_2,\dots, Q_{h}, F_2,&\dots, F_{h},\I_1,\dots, \I_{h},A,1/A)\\
	&= g(Q_2,\dots, Q_{h}, F_2,\dots, F_{h}, P_2(\I,F,Q),\dots, P_h(\I,F,Q),A,1/A)
\end{aligned}
\]
Let us now denote by $k_1$ the largest $h\geq 2$ such that $f$ depends non-trivially on $Q_h$; similarly let us denote by  $k_2$ the largest $h\geq 2$ such that $f$ depends non-trivially on $F_h$.
By definition $f$ is a constant of motion, so 
\[
0= \frac d{dt}f = \partial_{Q_{k_1}} f  \dot Q_{k_1} +
\partial_{F_{k_2}} f  \dot F_{k_2} + \mbox{other terms}\,.
\]
Now, if $k_2\geq k_1$ the only term in the expression above which contains $ Q_{k_2+1}$ is
$-4 A Q_{k_2+1}\partial_{F_{k_2}} f$, which in turn means that  $\partial_{F_{k_2}} f$ must be zero and $f$ cannot depend on $F_{k_2}$.
In the same way if $k_1>k_2$ then
\[
\frac d{dt}f=\partial_{Q_{k_1}} f A F_{k_1} + \mbox{other terms} =0
\]  
by construction all the other terms cannot contain $F_{k_1}$ so again we must have $\partial_{Q_{k_1}} f =0$ and hence a trivial dependence on $Q_{k_1}$.
This is only possible if $f$ does not depend on $F,Q$.
Now suppose that $f= P(\I,A,1/A)$ is a constant of motion. Passing to the common denominator, we have  that $P(A,1/A)=F(A)/A^M$, where  $F$  is a polynomial and $M$ a non-negative integer. Then
$$
\frac{d}{dt}P(A,1/A)= \frac{d}{dt}\left( \frac{F(A)}{A^M} \right) = \frac{\dot A}{A^{M+1}}\left(AF'(A)-MF(A)\right)=0,
$$
so that  $AF'(A)=MF(A)$, thus $F(A)=\alpha A^M$ and $P(A,1/A)$ does not depend on $A$.
This proves that any $g\in \mathscr{C}$ is in fact a finite polynomial in the $\I_k$.

\smallskip
Regarding the last statement in the Proposition,
the fact that all the $\I_k$ are in involution with $\I_1$, defined in \eqref{def:J1}, is the definition of constant of motion. 
\\
To prove that $\{\I_{k_1},\I_{k_2}\} =0$  with $k_i\geq 2$ we proceed by contradiction. Suppose that for some $k_1,k_2$ we have
$g(E,F,Q)= \{\I_{k_1},\I_{k_2}\} \not\equiv 0$. By the Poisson algebra property and Jacobi identity  we have that $g\in \mathscr{C}$ hence, as we proved above, we must have $g(E,F,Q)= f(\I(E,F,Q)) $ for some polynomial $f$. Since all the $\I_k$ contain only even powers of the variables $Q_i$ but, by Lemma \ref{involuz}, $g$ is linear in $Q$, we obtain a contradiction.   
\end{proof}

\section{Generating functions}\label{gege}
Once one has a family of constants of motion in involution, it is natural to associate to them a generating function. In this section we construct this function at a purely formal level, without discussing its domain of definition and its regularity. 
\begin{equation}\label{EFQlambda}
	D(\lambda):= \sum_{k=2}^\infty \lambda^{k-1} D_k, \quad B(\lambda):= \sum_{k=2}^\infty \lambda^{k-1} B_k\,,
	\quad Q(\lambda):= \sum_{k=2}^\infty \lambda^{k-1} Q_k.
\end{equation} and coherently
\[
\I(\lambda):=\sum_{k=2}^\infty \lambda^{k-1} \I_k, \quad E(\lambda):= \sum_{k=2}^\infty \lambda^{k-1} E_k, \quad F(\lambda):= \sum_{k=2}^\infty \lambda^{k-1} F_k\,.
\]
Note that the formula above corresponds to formula (3.5) of \cite{campos} (provided that $\lambda\to -z$).
\begin{lemma}\label{lemma generating function}
	The following identity holds:
	\begin{equation}\label{generating function}
		\begin{aligned}
			\I(\lambda)&=E(\lambda)-\frac{1}{\lambda}Q^2(\lambda)+\frac{1}{4}(E^2(\lambda)-F^2(\lambda)) \\ &= E(\lambda) - \frac{1}{\lambda}Q^2(\lambda) + D(\lambda)B(\lambda)= (D(\lambda)+ \frac1A)(B(\lambda) +A) -\frac1\lambda Q^2-1
		\end{aligned}
	\end{equation}
\end{lemma}
\begin{proof}
	By construction taking the terms of homogeneous degree $k-1$ on both sides of the equation.
	\[
	\I_k = E_k +\frac14 \sum_{\substack{k_1,k_2\geq 2\\ k_1+k_2 = k+1}} E_{k_1} E_{k_2} - F_{k_1} F_{k_2} - 
	\sum_{\substack{k_1,k_2\geq 2\\ k_1+k_2 = k+2} }Q_{k_1} Q_{k_2} 
	\]	
	which is exactly \eqref{def I}.
\end{proof}
Note that the building blocks of the generating function can be expressed at least formally in terms of pseudo-differential operators
\begin{equation}
	\label{cita}
	D(\lambda) = \sum_{j \neq 0} \frac{\lambda |j|^4}{1 - \lambda |j|^2} |u_j|^2\,,\quad B(\lambda) = \sum_{j \neq 0} \frac{\lambda |j|^2}{1 - \lambda |j|^2} |v_j|^2\,,\quad  Q(\lambda) = \sum_{j \neq 0} \frac{\lambda |j|^2}{1 - \lambda |j|^2}  u_j \bar v_j
\end{equation}
so in conclusion
%\[
%\begin{aligned}
%\I(\lambda) = \, &\lambda \left( 1 + \|\nabla u\|_{L^2}^2 \right) \left\| \frac{\nabla}{\sqrt{1 + \lambda \Delta}} \, v \right\|_{L^2}^2 + \left( \frac{\lambda}{1 + \|\nabla u\|_{L^2}^2} \right) \left\| \frac{\Delta}{\sqrt{1 + \lambda \Delta}} \, u \right\|_{L^2}^2 \\[15pt]
%	&- \lambda \left( \left\langle \frac{1}{1 + \lambda \Delta} \, \nabla u, \, \nabla v \right\rangle_{L^2} \right)^2 + \lambda^2 \left\| \frac{\Delta}{\sqrt{1 + \lambda \Delta}} \, u \right\|_{L^2}^2 \left\| \frac{\nabla}{\sqrt{1 + \lambda \Delta}} \, v \right\|_{L^2}^2\,.
%\end{aligned}
%\]
%\red{ provo un altra scrittura}
\[
\I(\lambda)= (\sum_{j\ne 0}  \frac{ |j|^2}{1 - \lambda |j|^2} |u_j|^2+1)(\sum_{j \neq 0}  \frac{\lambda |j|^2}{1 - \lambda |j|^2} |v_j|^2 +A) -\frac1\lambda (\sum_{j \neq 0} \frac{\lambda |j|^2}{1 - \lambda |j|^2} u_j \bar v_j)^2 -1\,.
\]
It may be convenient to set $\mathcal J(\lambda):= \I(\lambda)-\I_1$, where $\I_1=A-B_1-1$ is defined in \eqref{def:J1},  so that 
\[
\begin{aligned}
	\mathcal J(\lambda)&=\left( \sum_{j \neq 0} \frac{|j|^2}{1 - \lambda |j|^2} |u_j|^2 \right) \left( \sum_{j \neq 0} \frac{\lambda |j|^2}{1 - \lambda |j|^2} |v_j|^2 + A \right) + \sum_{j \neq 0} \frac{1}{1 - \lambda |j|^2} |v_j|^2 - \frac{1}{\lambda} \left( \sum_{j \neq 0} \frac{\lambda |j|^2}{1 - \lambda |j|^2} u_j \bar{v}_j \right)^2\\
	=& \sum_{j \neq 0}  \frac{ A |j|^2  |u_j|^2+|v_j|^2}{1 - \lambda |j|^2} + \lambda \left[\left( \sum_{j \neq 0} \frac{|j|^2}{1 - \lambda |j|^2} |u_j|^2 \right) \left( \sum_{j \neq 0} \frac{ |j|^2}{1 - \lambda |j|^2} |v_j|^2 \right) - \left( \sum_{j \neq 0} \frac{ |j|^2}{1 - \lambda |j|^2} u_j \bar{v}_j \right)^2\right]
	\\= & \sum_{j \neq 0}  \frac{ A |j|^2  |u_j|^2+|v_j|^2}{1 - \lambda |j|^2} + \frac{\lambda}{2} \sum_{j \neq 0} \sum_{k \neq 0} \frac{|j|^2 |k|^2}{(1 - \lambda |j|^2)(1 - \lambda |k|^2)} |u_j v_k - u_k v_j|^2\\
	&= \sum_{j \neq 0} \frac{1}{1 - \lambda |j|^2} \left[ A |j|^2 |u_j|^2 + |v_j|^2 + \sum_{\substack{k \neq 0 \\ |k|^2 \neq |j|^2}} \frac{|j|^2 |k|^2}{|j|^2 - |k|^2} |u_j v_k - u_k v_j|^2 \right] \\
	&\quad + \frac{\lambda}{2} \sum_{\substack{j, k \neq 0 \\ |j|^2 = |k|^2}} \frac{|j|^4}{(1 - \lambda |j|^2)^2} |u_j v_k - u_k v_j|^2=\sum_{\nu \in \Lambda} \left[ \frac{1}{1 - \lambda \nu} \mathcal I_\nu^{(1)} + \frac{\lambda \nu^2}{(1 - \lambda \nu)^2} \mathcal I_\nu^{(2)} \right]
\end{aligned}
\]
Taylor expanding at $\lambda=0$ we obtain at least at a formal level the relations \eqref{calIJ}.
%A direct computation ensures that
%\begin{equation}
%	\begin{aligned}
	%		\mathcal J(\lambda) &= \sum_{j \neq 0} \frac{1}{1 - \lambda |j|^2} \left[ A |j|^2 |u_j|^2 + |v_j|^2 + \sum_{\substack{k \neq 0 \\ |k|^2 \neq |j|^2}} \frac{|j|^2 |k|^2}{|j|^2 - |k|^2} |u_j v_k - u_k v_j|^2 \right] \\
	%		&\quad + \frac{\lambda}{2} \sum_{\substack{j, k \neq 0 \\ |j|^2 = |k|^2}} \frac{|j|^4}{(1 - \lambda |j|^2)^2} |u_j v_k - u_k v_j|^2=\sum_{\nu \in \Lambda} \left[ \frac{1}{1 - \lambda \nu} \mathcal I_\nu^{(1)} + \frac{\lambda \nu^2}{(1 - \lambda \nu)^2} \mathcal I_\nu^{(2)} \right]
	%	\end{aligned}
%\end{equation}
We have formally recovered the constants of motion $\cI_\nu^{(i)}$ from the generating function.
%
%So that
%\begin{equation}
%	J(\lambda) = \sum_{\nu \in \Lambda} \left[ \frac{1}{1 - \lambda \nu} \mathcal I_\nu^{(1)} + \frac{\lambda \nu^2}{(1 - \lambda \nu)^2} \mathcal I_\nu^{(2)} \right]
%\end{equation}
%Taylor expanding in $\lambda=0$ with
%\[
%\frac{1}{1-\lambda\nu} \sim \sum_{h=0}^\infty (\lambda \nu)^h \,,\qquad \frac{\lambda \nu^2}{(1 - \lambda \nu)^2} \sim \sum_{m=1}^\infty m \lambda^m \nu^{m+1}
%\]
%\[
%J(\lambda)= -\I_1+\sum_{k\geq 2} \I_k\lambda^{k-1} = \sum_{\nu \in \Lambda}  \sum_{k=1}^\infty \lambda^{k-1} ( \nu^{k-1} \mathcal I^{(1)}_\nu + (k-1) \nu^{k} \mathcal I^{(2)}_\nu)
%\]
%so that
%\[
%\I_1=  -\sum_{\nu \in \Lambda}\mathcal I^{(1)}_\nu \,,\quad \I_k = \sum_{\nu \in \Lambda} \nu^{k-1} \mathcal I^{(1)}_\nu + (k-1) \nu^{k} \mathcal I^{(2)}_\nu
%\]
%o se ci piace di piu'
%\[
%\I(\lambda)= (\sum_{j\ne 0}  \frac{ |j|^2}{1 - \lambda |j|^2} |u_j|^2+1)(\sum_{j \neq 0}  \frac{1}{1 - \lambda |j|^2} |v_j|^2 +\I_1) -\frac1\lambda (\sum_{j \neq 0} \frac{\lambda |j|^2}{1 - \lambda |j|^2} u_j \bar v_j)^2-1
%\]
%\red{DARIO GUARDA QUI, please:} in entrambe i casi si ottiene la prima formula di Dario (il primissimo file con la matrice due per due):
%\begin{equation}\label{davedere}
%\I(\lambda)+1 = \det \mathcal L\,,\quad \mathcal L= \begin{pmatrix} D(\lambda)+ 1/A & \frac{Q(\lambda)}{\sqrt{\lambda}} \\
	%\frac{Q(\lambda)}{\sqrt{\lambda}} & B(\lambda)+ A \end{pmatrix}
%\end{equation}
% dove (almeno a livello formale, sto ignorando ogni problema legato alla risommazione che ho fatto il \ref{cita} e \ref{cita2}...)
% \[
% D(\lambda)+ 1/A= D(\lambda)+D_1 +1 = \sum_{j\ne 0}  \frac{ |j|^2}{1 - \lambda |j|^2} |u_j|^2+1\,,
% \]
% \[
%  B(\lambda)+ A = B(\lambda)+ B_1 +(A-B_1)=  \frac{1}{1 - \lambda |j|^2} |v_j|^2 +\I_1
% \]
% \begin{equation}
	% 	\I(\lambda) - \I_1 +1 = \left( \sum_{j \neq 0} \frac{|j|^2}{1 - \lambda |j|^2} |u_j|^2 \right) \left( \sum_{j \neq 0} \frac{\lambda |j|^2}{1 - \lambda |j|^2} |v_j|^2 + A \right) + \sum_{j \neq 0} \frac{1}{1 - \lambda |j|^2} |v_j|^2 - \frac{1}{\lambda} \left( \sum_{j \neq 0} \frac{\lambda |j|^2}{1 - \lambda |j|^2} u_j \bar{v}_j \right)^2
	% \end{equation}
% \begin{align}
	% 	\I(\lambda) +1 - \I_1 &= \left[ \sum_{j \neq 0} \frac{|j|}{1 - \lambda |j|^2} \Big( |z_j|^2 + \text{Re}(z_j z_{-j}) \Big) \right] \left[ \lambda \sum_{k \neq 0} \frac{|k|^3}{1 - \lambda |k|^2} \Big( |z_k|^2 - \text{Re}(z_k z_{-k}) \Big) \right] \nonumber \\
	% 	&\quad + \left[ \sum_{j \neq 0} \frac{|j|}{1 - \lambda |j|^2} \Big( |z_j|^2 + \text{Re}(z_j z_{-j}) \Big) \right] \frac{1}{1 + \sum_{k \neq 0} |k| \Big( |z_k|^2 + \text{Re}(z_k z_{-k}) \Big)} \nonumber \\
	% 	&\quad + \sum_{j \neq 0} \frac{|j|}{1 - \lambda |j|^2} \Big( |z_j|^2 - \text{Re}(z_j z_{-j}) \Big) \nonumber \\
	% 	&\quad - \lambda \left[ \sum_{j \neq 0} \frac{|j|^2}{1 - \lambda |j|^2} \text{Im}(z_j z_{-j}) \right]^2
	% \end{align}
% 	Let us consider the symplectic change of variables
% \begin{align}
	% 	u_j &= \frac{1}{\sqrt{2|j|}} \left( z_j + \bar{z}_{-j} \right) \\
	% 	v_j &= \sqrt{\frac{|j|}{2}} \frac{z_j - \bar{z}_{-j}}{i}
	% \end{align}
% in these variables we obtain
% \begin{equation}
	% 	H = (U + 1)(V + A) - \frac{1}{\lambda} M^2
	% \end{equation}
% where
% \begin{align}
	% 	U &= \sum_{j \neq 0} \frac{|j|}{1 - \lambda |j|^2} \Big( |z_j|^2 + \text{Re}(z_j z_{-j}) \Big) \\
	% 	V &= \lambda \sum_{j \neq 0} \frac{|j|^3}{1 - \lambda |j|^2} \Big( |z_j|^2 - \text{Re}(z_j z_{-j}) \Big) \\
	% 	M &= \lambda \sum_{j \neq 0} \frac{|j|^2}{1 - \lambda |j|^2} \text{Im}(z_j z_{-j}) \\
	% 	X &= \sum_{j \neq 0} |j| \Big( |z_j|^2 + \text{Re}(z_j z_{-j}) \Big) \\
	% 	A &= \frac{1}{1 + X}
	% \end{align}
% whose quadtatic part is
% \begin{equation}
	% 	H_2 = \sum_{j \neq 0} \frac{2 \lambda |j|^3}{1 - \lambda |j|^2} |z_j|^2
	% \end{equation}
% \red{Ecco la domanda:} se mi metto su un supporto finito per esempio $u_j=v_j=0$ per ogni $j\notin \{j_1,\dots,j_N\}$. Questa formula coincide con la tua formula (24)? vedo ingredienti simili ma faccio fatica a mettere in relazione le due cose.

\medskip



%\begin{remark}
%	While we preferred working on tori-- both  for simplicity and in order to show the formal integrability in the $1d$-case-- the results in this section and the next one are quite general and only require to be able to integrate by parts.
%\end{remark}

\subsection{Expansion of the generating function at infinity.}
As we have seen in formula  \ref{cita}, we may write at least at the formal level
\[
D(\lambda) = \sum_{j \neq 0} \frac{\lambda |j|^4}{1 - \lambda |j|^2} |u_j|^2 \,,\quad B(\lambda) = \sum_{j \neq 0} \frac{\lambda |j|^2}{1 - \lambda |j|^2} |v_j|^2 \,,\quad Q(\lambda) = \sum_{j \neq 0} \frac{\lambda |j|^2}{1 - \lambda |j|^2}{ u_j \bar v_j }
\]
so we may apply the formal change of variables $\Phi:\lambda\to \kappa= 1/\lambda$, we set
\[
\tilde D(\kappa) = -D(\frac1\kappa) = \sum_{j \neq 0} \frac{ |j|^2}{1 - \kappa |j|^{-2}} |u_j|^2 \,,\quad \tilde B(\kappa) =  \sum_{j \neq 0} \frac{ 1}{1 - \kappa |j|^{-2}} |v_j|^2 \,,\quad  \tilde Q(\kappa) =  \sum_{j \neq 0} \frac{ 1}{1 - \kappa |j|^{-2}} u_j\bar v_j 
\]
so that
\[
\tilde \I(\kappa)=\I(\frac1\kappa) = - A \tilde D(\kappa) -\frac{\tilde B(\kappa)}{A} + \tilde D(\kappa) \tilde B(\kappa) -\kappa \tilde Q(\kappa)^2
\]
Now we Taylor expand at $\kappa=0$ to obtain constants of motion involving the negative Sobolev norms.
At the formal level we have 
\[
\hat D(\kappa) =\sum_{h=-1}^\infty \kappa^{h+1} D_{-h}\,, \hat B(\kappa) =\sum_{h=-1}^\infty \kappa^{h+1} B_{-h}\,,\quad \hat Q(\kappa) =\sum_{h=-1}^\infty \kappa^{h+1} Q_{-h}
\]
where we have extended in the trivial way the building blocks $D_k,B_k,Q_k$  defined in \ref{BDQ} to negative values of $k$.
Then, for any $h\geq -1$ one has that
\begin{equation}\label{Jnegative}
	\I_{-h}:= - A  D_{-h} -\frac{ B_{-h}}{A} + \sum_{\substack{h_1,h_2\geq -1\\ h_1+h_2= h-1}}  D_{-h_1}B_{-h_2} -
	\sum_{\substack{h_1,h_2\geq -1\\ h_1+h_2= h-2}} Q_{-h_1} Q_{-h_2} 
\end{equation}
are constants of motion.\\
In the case $h>-1$, by expanding $B_{-h}/A$ as $B_{-h}+D_1B_{-h}$ and the second summation in \eqref{Jnegative} as
\begin{equation*}
	\sum_{\substack{h_1,h_2\geq -1\\ h_1+h_2= h-1}}  D_{-h_1}B_{-h_2}=\sum_{\substack{h_1,h_2\geq 0\\ h_1+h_2= h-1}}  D_{-h_1}B_{-h_2}+
	D_1B_{-h}+D_{-h}B_1
\end{equation*}
we can write
\begin{equation*}
	\I_{-h}= (B_1-A)D_{-h} -B_{-h} + \sum_{\substack{h_1,h_2\geq 0\\ h_1+h_2= h-1}}  D_{-h_1}B_{-h_2} -
	\sum_{\substack{h_1,h_2\geq -1\\ h_1+h_2= h-2}} Q_{-h_1} Q_{-h_2}\\
	%&=-\I_1D_{-h}-D_{-h} -B_{-h} + \sum_{\substack{h_1,h_2\geq 0\\ h_1+h_2= h-1}}  D_{-h_1}B_{-h_2} -
	%\sum_{\substack{h_1,h_2\geq -1\\ h_1+h_2= h-2}} Q_{-h_1} Q_{-h_2}
\end{equation*} 
Explicilty
\begin{align*}
	\I_{1} &= A - B_{1} - 1 \,,\quad 	\I_{0}= (B_1-A)D_{0} - B_{0} - Q_{1}^2 \,,\quad 
	\I_{-1} = (B_1-A)D_{-1}- B_{-1}+ D_{0}B_{0}  - 2Q_{1}Q_{0} \\
	\I_{-2} &= (B_1-A)D_{-2}- B_{-2}+ D_{0}B_{-1} + D_{-1}B_{0}  - 2Q_{1}Q_{-1} - Q_{0}^2\\
	\I_{-3} &= (B_1-A)D_{-3}- B_{-3} + D_{0}B_{-2} + D_{-1}B_{-1} + D_{-2}B_{0}  - 2Q_{1}Q_{-2} - 2Q_{0}Q_{-1} \\
	\I_{-4} &= (B_1-A)D_{-4}- B_{-4}  + D_{0}B_{-3} + D_{-1}B_{-2} + D_{-2}B_{-1} + D_{-3}B_{0}  - 2Q_{1}Q_{-3} - 2Q_{0}Q_{-2} - Q_{-1}^2\,.
\end{align*}
Note that $\I_1$ is the one defined in \eqref{def:J1}.
Recalling the definition of $\mathcal J(\lambda)$ we note that, at least formally
\[
\mathcal J(\frac1\kappa)= %\sum_{\nu \in \Lambda} \left( \frac{\kappa}{\kappa -  \nu} \mathcal I_\nu^{(1)} + \frac{\kappa\nu^2}{ (\kappa -  \nu)^2} \mathcal I_\nu^{(2)} \right)=
\sum_{\nu \in \Lambda} \left( \frac{\kappa \nu^{-1}}{\kappa\nu^{-1} - 1} \mathcal I_\nu^{(1)} + \frac{\kappa}{ (\kappa\nu^{-1} -  1)^2} \mathcal I_\nu^{(2)} \right)= \sum_{h=0}^\infty \kappa^{h+1} \I_{-h}
\]
Taylor expanding at $\kappa=0$ we obtain at least at a formal level
\begin{equation}\label{Inegative}
	\I_{-h}= \sum_{\nu\in \Lambda} \Big(- \nu^{-h-1}\cI_\nu^{(1)} + (h+1) \nu^{-h} \cI^{(2)}_\nu\Big)\,.
\end{equation}

\section{Constants of motion in the energy space}
The aim of this section is to prove Theorem \ref{genero}. We proceed in various step proving first that the $\mathcal I^{(n)}_\nu$ are well defined, then that they generate the $\I_k$ and finally that they are in involution. While these results could be proved by direct computations, our approach relies on the density of trigonometric polynomials in $\H^s$. Note that, on Galerkin truncations, the functions $\cI$ are polynomials, the  function $\mathcal{J}(\lambda)$ is rational and hence the formal identities \eqref{calIJ} and \eqref{Inegative} are true identities.
\begin{lemma}\label{tameI}
If $(u, v) \in \H^1$, then for every $\nu \in \Lambda$, the functions  $\mathcal I_\nu^{(n)}$ are well-defined, and the infinite series in \eqref{eq:Gamma1} converges absolutely. Finally 
\begin{equation}
	\label{uno}
	\sum_{\nu\in \Lambda}  \mathcal I_\nu^{(1)} \leq 	\|(u,v)\|^2_{\H^{1}}+ C \|(u,v)\|^2_{\H^{1}} \|(u,v)\|^2_{\H^{2}}\,,
\end{equation}
while, for $k>0$ we have the tame estimate
\begin{equation}
	\label{due}
	\sum_{\nu\in \Lambda} \nu^k \mathcal I_\nu^{(1)} \leq 	\|(u,v)\|^2_{\H^{k+1}}+ C \|(u,v)\|^2_{\H^{k+1}} \|(u,v)\|^2_{\H^{2}}\,.
\end{equation}
	\end{lemma}
	
	\begin{proof}
Fix an energy level $\nu \in \Lambda$. 
The fact that the $M_j$ and $\mathcal I_\nu^{(2)}$ are well defined is obvious, since they are finite sums.  Moreover, by definition we have that
$0<A\leq \frac{1}{1+\|u\|_{H^1}^2}\leq 1$.
\\
It remains to prove the absolute convergence of the infinite sum in \eqref{eq:Gamma1} for any fixed $j \in S_\nu$:
\begin{equation}\label{eq:Aj}
	A_j = \sum_{\Lambda\ni \nu'\ne \nu} \sum_{k \in S_{\nu'}} \left| \frac{\nu \nu' }{\nu - \nu'} \right| |u_j v_k - u_k v_j|^2.
\end{equation}
recalling that $ |u_j v_k - u_k v_j|^2 \leq 2\left( |u_j|^2 |v_k|^2 + |v_j|^2 |u_k|^2 \right) $, that $|\nu-\nu'|\geq 1$ and 
\begin{equation}\label{piacimento}
	\sup_{\nu' \in \Lambda \setminus \{\nu\}} \left| \frac{\nu \nu'}{\nu - \nu'} \right| \leq \nu(\nu+1)
\end{equation} we obtain 
\[
|A_j |\leq 2 \nu( \nu |u_j|^2   \sum_{k\notin S_\nu} |v_k|^2 +  |v_j|^2 \sum_{k\notin S_\nu}|k|^2 |u_k|^2) \leq 2 \nu^2 |u_j|^2  \|v\|_{L^2}^2 + 2\nu |v_j|^2 \|u\|_{H^1}^2
\]
which imples \eqref{uno}
To prove \eqref{due}, consider 
\[
\begin{aligned}
	\sum_{j\in \Z^n} |j|^{2 k} A_j &= \sum_{\nu} \sum_{\nu' \neq \nu} \nu^k \left\vert{} \frac{\nu \nu'}{\nu - \nu'} \right\vert{} \sum_{j \in S_\nu} \sum_{k \in S_{\nu'}} \vert{}u_j v_k - u_k v_j\vert{}^2 \\ &\leq  2 \sum_{\nu \neq \nu'} \nu^k \left\vert{} \frac{\nu \nu'}{\nu - \nu'} \right\vert{} \Big( U_\nu V_{\nu'} + V_\nu U_{\nu'} \Big)\\ & \leq 2\sum_{\nu \neq \nu'}  \left\vert{} \frac{\nu \nu'}{\nu - \nu'} \right\vert{}(\nu^k+(\nu')^k) U_\nu V_{\nu'} 
\end{aligned}
\]
where we have set for compactness of notation
\begin{equation}
	\label{blocchi}
	U_\nu:=\sum_{j\in S_\nu} |u_j|^2\,,\quad 	V_\nu:=\sum_{j\in S_\nu} |v_j|^2\,.
\end{equation}
Now we bound (using also \eqref{piacimento})
\begin{equation}\label{estimatetame}
	\frac{\nu \nu'}{|\nu - \nu'|} (\nu^k+(\nu')^k)  \leq	\begin{cases}
		2(\nu')^{k} \nu(\nu+1) \qquad &\mbox{if}\quad \nu'>\nu
		\\ 2(\nu)^{k+1} \nu' \qquad &\mbox{if}\quad \nu'<\nu
	\end{cases}
\end{equation}
which implies \eqref{due}.
\end{proof}

%If one is only interested in convergence then some interesting cancellations may occur 
%and one can obtain the bound
%$$\sum_{\nu \in \Lambda} \nu^k \cI_\nu^{(1)} \leq A \Vert{}\nabla^{k+1} u\Vert{}_{L^2}^2 + \Vert{}\nabla^k v\Vert{}_{L^2}^2 + 2 \sum_{m=0}^{k-1} \Vert{}\nabla^{m+1} u\Vert{}_{L^2}^2 \Vert{}\nabla^{k-m} v\Vert{}_{L^2}^2$$ as can be seen in the following result.
\begin{lemma}\label{Iinverso}
Let $\I_k$, $\mathcal I^{(1)}_\nu$, and $\mathcal I^{(2)}_\nu$ be defined as above. Then, the following identities hold:
\begin{equation}\label{Iuno}
	\I_1 = -\sum_{\nu \in \Lambda} \mathcal I^{(1)}_\nu
\end{equation}
and for all $k \geq 2$,
\begin{equation}\label{Ik}
	\I_k = \sum_{\nu \in \Lambda} \nu^{k-1} \mathcal I^{(1)}_\nu + (k-1) \sum_{\nu \in \Lambda} \nu^{k} \mathcal I^{(2)}_\nu\,,
\end{equation}
where the right-hand side of \eqref{Iuno} and \eqref{Ik} are absolutely convergent series respectively in $\H^{2}$ and $\H^{k}$.
\end{lemma}

\begin{proof}
	For every $N\in\N_{\geq 1}$, let $\Lambda_N$ be the set, $\{\nu_1,\ldots,\nu_N\}$, of the first $N$ elements of the ordered set $\Lambda$ and set $\mathscr{U}_N:=\{ (u,v)\in \H^1\,\vert{} \quad u_j=v_j=0\,,\quad \forall |j|^2>\nu_N\}$ and consider the functions $\I_k$ and the  $\cI^{(i)}_{\nu_1},\ldots,\cI^{(i)}_{\nu_N}$, $i=1,2$ restricted on $\mathscr{U}_N$.\\
	Under this restriction,  the $\I,\cI$ are polynomials in a finite number of variables while $\mathcal{J}(\lambda)$ is a rational function in $\lambda$ which is analytic at zero. Then, Taylor expanding as explained in Section \ref{gege} we obtain the result as a true identity (and not just a formal one).
	Since the spaces $\mathscr{U}_N$ are dense in $\H^k$ in order to conlcude our proof we only need to show that both the left and the right hand side of \eqref{Iuno} and \eqref{Ik} are bounded polynomials and hence (see \cite{muj})  analytic functions. The fact that the $\I_k$ are bounded polynomials  in $\H^k$ follows directly from their definition. The fact that the right hand sides of  \eqref{Iuno} and \eqref{Ik} are formal polynomials comes form the fact that \eqref{eq:Aj} is formally well defined. It remains to prove the boundedness, respectively in $\H^2$ and $\H^k$.  The boundedness of the right hand side of \eqref{Iuno} follows from \eqref{uno}. Regarding \eqref{Ik} the boundedness of the first summand follows from \eqref{due} with $k\rightsquigarrow k-1$.  The boundedness of the second  summand follows from 
	\[
	\sum_{\nu\in\Lambda}\nu^k \cI^{(2)}_\nu \leq \|u\|_{H^k}^2\|v\|_{L^2}^2\,.
	\]
	
\end{proof}
	
	
	
	
	
	

\begin{lemma}\label{prop:comm_gamma2}
On the space $\H^3$, the quantities $\cI_\nu^{(1)}$, $\cI_\nu^{(2)}$, and $M_j$ form a set of mutually commuting functions with respect to the Darboux Poisson bracket. That is, for any indices $j,k\in \Z^n$ and $\nu, \mu\in \Lambda$:
\begin{enumerate}
	\item $\{M_j, M_k\} = \{\cI_\nu^{(2)}, M_k\} = \{\cI_\nu^{(2)}, \cI_\mu^{(2)}\} =\{\cI_\nu^{(1)}, M_k\} = 0$,
	\item $ \{\cI_\nu^{(1)}, \cI_\mu^{(2)}\} = 0$,
	\item $\{\cI_\nu^{(1)}, \cI_\mu^{(1)}\} = 0$.
\end{enumerate}
\end{lemma}

\begin{proof}
The fact that the $\cI^{(2)}_\nu,M_j$ are all in involution is trivial, as well as the fact that $\cI^{(1)}_\nu$ Poisson commutes with all the $M_j$. Therefore item (1) follows.\\
In order to deal with the remaining brackets we divide the proof in three steps:\\
{\bf Step 1. } We first prove the item (2)  restricted to $\mathscr{U}_N$, so that all the sums are finite.
It is useful to express  the $\cI^{(n)}_\nu$ in terms of the blocks $U,V,W$ where $W_\nu:=\sum_{j\in S_\nu} u_j\bar v_j$\,.
We have
\begin{equation}\label{eq:Gamma1_macro}
	\cI_\nu^{(1)} = A \nu U_\nu + V_\nu + \sum_{\nu' \ne \nu} \frac{\nu \nu'}{\nu - \nu'} \big( U_\nu V_{\nu'} + V_\nu U_{\nu'} - 2 W_\nu W_{\nu'} \big),\quad \cI_{\nu}^{(2)} = U_{\nu}V_{\nu} - W_{\nu}^2
\end{equation}
using the commutation rules we have
\begin{equation}
	\{U_\nu, V_\nu\} = 4W_\nu, \quad \{U_\nu, W_\nu\} = 2U_\nu, \quad \{V_\nu, W_\nu\} = -2V_\nu
\end{equation}
(of course trivially $\{U_\nu, U_\mu\} = 0$, etc., for $\nu \ne \mu$).
From the above commutation rules, one can easily verify that $\cI^{(2)}_\nu$ commutes with the quantities $U_\mu,\,V_\mu,\,W_\mu$ for all $\nu,\mu\in\Lambda$, so that, item (2) follows directly. Moreover, by \eqref{Ik} together with item (2) , we have that $\cI_\nu^{(2)}$ also commutes with all the $\I_k$.\\
{\bf Step 2. }We study item (3) restricted to $\mathscr{U}_N$. Thanks to Lemma \ref{Iinverso} we have the linear system
\begin{equation}
	\begin{pmatrix}
		-\I_1\\
\J_2\\		\vdots\\
		\I_N
	\end{pmatrix}=\M_1
	\begin{pmatrix}
		\cI^{(1)}_{\nu_1}\\
			\cI^{(1)}_{\nu_2}\\
		\vdots\\
		\cI^{(1)}_{\nu_N}
	\end{pmatrix}+\M_2\begin{pmatrix}
		\cI^{(2)}_{\nu_1}\\
			\cI^{(2)}_{\nu_2}\\
		\vdots\\
		\cI^{(2)}_{\nu_N}
	\end{pmatrix},
\end{equation}
where
\begin{equation}
	\M_1=\begin{pmatrix}
		1&\ldots&1\\
		\nu_1&\ldots&\nu_N\\
		\vdots&&\vdots\\
		\nu_1^{N-1}&\ldots&\nu_N^{N-1}
	\end{pmatrix},\quad \M_2=\begin{pmatrix}
		0&\ldots&0\\
		\vdots&&\vdots\\
		(k-1)\nu_1^k&\ldots&(k-1)\nu_N^k\\
		\vdots&&\vdots\\
		(N-1)\nu_1^N&\ldots&(N-1)\nu_N^N
	\end{pmatrix}.
\end{equation}
Since $\I$ and $\cI^{(2)}$ are in involution and $\M_1$ is a Vandermonde matrix and hence invertible,  we have that item (3) restricted to $\mathscr{U}_N$ holds for any $N\in\N$.\\

{\bf Step 3.}  It remains to show that the expressions in items (1),(2),(3) are bounded polynomials in $\H^3$. We only discuss the case of item (3), which is the most complicated.
 Since momentum preserving formal power series are closed with respect to formal Poisson brackets (see for instance \cite{ProcesiStolo}) we have that the Poisson bracket between $\cI^{(1)}_\nu$ and $\cI^{(1)}_\mu$ is a formal polynomial.
 In performing the bracket $\{\cI_\nu^{(1)}, \cI_\mu^{(1)}\} $  several terms arise. Let us consider explicitly
 \[
\begin{aligned}
	&\nu\mu  U_\nu  \sum_{\mu' \ne \mu}  \frac{\mu'}{\mu - \mu'} \{	A ,( U_\mu V_{\mu'} + V_\mu U_{\mu'} - 2 W_\mu W_{\mu'}) \}
\\
&= -A^2 \nu\mu  U_\nu  \sum_{\mu' \ne \mu}  \frac{\mu'}{\mu - \mu'} \{	\sum_{\nu'} \nu' U_{\nu'} ,( U_\mu V_{\mu'} + V_\mu U_{\mu'} - 2 W_\mu W_{\mu'}) \}
\end{aligned}
 \]
 There are two cases, either $\nu'=\mu$ or $\nu'=\mu'$.  If $\nu'=\mu$ the worst case scenario is a term of the type $\sum_{\mu' \ne \mu}\mu' V_{\mu'}$ which is convergent provided that $(u,v)\in \H^2$. If $\nu'=\mu'$ then  the worst case scenario is a term of the type $\sum_{\mu' \ne \mu}(\mu' )^2V_{\mu'}$ which is convergent provided that $(u,v)\in \H^3$.  All the other terms can be estimated similarly.
 
% 
%Let us denote by $P_{\nu,\mu}$ the Poisson between $\cI^{(1)}_\nu$ and $\cI^{(1)}_\mu$. Since its restriction, $P_{\nu,\mu, N}$, on $\mathscr{U}_N$ coincides with the Poisson bracket of the restrictions of $\cI^{(1)}_\nu$ and $\cI^{(1)}_\mu$, we have the desired result recalling that $\lim_{N\to\infty}P_{\nu,\mu, N}=P_{\nu,\mu}$.

%	By construction, every modified energy $\I_m$ depends on the fields $u$ and $v$ exclusively through the macroscopic quantities $D_l$, $B_l$, and $Q_l$ (for various orders $l$). These macroscopic quantities are sums over eigenspaces $S_l$ of the microscopic blocks:
%	\[
%	Ul = \sum_{j \in S_l} u_j u_{-j}, \quad V_l = \sum_{j \in S_l} v_j v_{-j}, \quad W_l = \sum_{j \in S_l} u_j v_{-j}.
%	\]
%	Therefore, any functional $\I_m$ is invariant under orthogonal (or unitary) transformations within each eigenspace $S_l$. 
%	
%	The quantity $\mathcal I_\nu^{(2)}$ is built from the microscopic building blocks $W_{j,k} := u_j v_k - u_k v_j$ for $j, k \in S_\nu$, which represent the generators of rotations (angular momentum components) within the eigenspace $S_\nu$. 
%	
%	To show that $\{\I_m, \mathcal I_\nu^{(2)}\} = 0$, it suffices to check that any microscopic block ($Ul, V_l, W_l$) commutes under the Poisson bracket with $W_{j,k}$ for $j, k \in S_\nu$. Using the canonical brackets $\{u_a, v_b\} = \delta_{a,-b}$ and $\{u_a, u_b\} = \{v_a, v_b\} = 0$, we compute the bracket with $Ul$:
%	\[
%	\{Ul, u_j v_k\} = \left\{ \sum_{a \in S_l} u_a u_{-a}, u_j v_k \right\} = \sum_{a \in S_l} \left( u_a \{u_{-a}, v_k\} u_j + \{u_a, v_k\} u_{-a} u_j \right).
%	\]
%	Since $\{u_{-a}, v_k\} \neq 0$ only when $-a = -k$ (i.e., $a = k$), and $\{u_a, v_k\} \neq 0$ only when $a = -k$, the antisymmetrized combination vanishes identically:
%	\[
%	\{Ul, W_{j,k}\} = \{Ul, u_j v_k - u_k v_j\} = 0.
%	\]
%	Analogously, one checks that:
%	\[
%	\{V_l, W_{j,k}\} = 0 \quad \text{and} \quad \{W_l, W_{j,k}\} = 0.
%	\]
%	Since every macroscopic quantity $D_l, B_l, Q_l$ commutes with $W_{j,k}$ and its complex conjugate, it follows that any modified energy $\I_m$ commutes with every individual term $W_{j,k}$. 
%	
%	Finally, because $\mathcal I_\nu^{(2)}$ is a sum of squares and products of these terms $W_{j,k}$ (and their complex conjugates), the Leibniz rule for Poisson brackets guarantees that:
%	\[
%	\{\I_m, \mathcal I_\nu^{(2)}\} = 0 \quad \text{for all } m \geq 1 \text{ and } \nu \in \Lambda.
%	\]
%	This completes the proof. $\blacksquare$
\end{proof}
We conclude by proving that also the identities \eqref{Inegative} hold not only at formal level.
\begin{lemma}\label{invonegative}
	For every $h\geq -1$ the identity \eqref{Inegative} holds on
	on $\mathcal{H}^2$.
As a consequence, for every $h,k\in \Z$  one has
	\begin{equation*}\label{PoissonMista}
	\{\I_{h},\I_k\}=0 \qquad\text{ on }\mathcal{H}^{m}.
	\end{equation*}
	where $m= \max(2, h+1,k+1)$.
\end{lemma}
\begin{proof}
	We recall that  identity \eqref{Inegative} restricted to
	$\mathscr{U}_N$ holds not just formally. The left hand side of \eqref{Inegative} is well defined in $\H^2$ by \eqref{Jnegative}, and the right hand side is absolutely convergent in $\H^2$, reasoning as in Lemma \ref{tameI}. The first result follows.
	\\
	To prove the involution property we only need to prove that $	\{\I_{h},\I_k\}$ is a bounded polynomial on $\mathcal{H}^{m}$, this is done just as in Step 3 of the proof of Lemma \ref{prop:comm_gamma2}.
%	We only prove \eqref{PoissonNegativa}, while the second identity follows by the same argument:\\
%	We first show that $\mathcal{P}_{h_1,h_2}:=\{\I_{-h_1},\I_{-h_2}\}$ is a well
%	defined, continuous polynomial on $\H^1$.\\
%	By \eqref{Jnegative}, for a fixed $h\geq-1$, the only blocks of order $1$ appearing in $\I_{-h}$, are the $Q_1$ and the combination
%	$B_1-A$.\\
%	Then, by Lemma \eqref{relations1} (extended to all integer indices) the only bracket that generates a block of order $2$ is $\{A,B_1\}=-4A^2Q_2$, which occurs only within $\{B_1-A,B_1-A\}=0$. Finally, all the remaining brackets produce blocks of order $\leq1$. Then, e conclude that $\mathcal{P}_{h_1,h_2}$ is a bounded (hence analytic) polynomial on $\H^1$.\\
%	We now argue as in item (3) of \ref{prop:comm_gamma2}:\\
%	For every $N\in\N$, let
%	$\Lambda_N$ be the set $\{\nu_1,\ldots,\nu_N\}$ of the first $N$ elements of the	ordered set $\Lambda$ and set
%	$\mathscr{U}_N:=\{(u,v)\in\H^1 \,\vert\, u_j=v_j=0\,,\ \forall |j|^2>\nu_N\}$.
%	Since $U_\nu=V_\nu=W_\nu=0$ on $\mathscr{U}_N$ for $\nu\notin\Lambda_N$, the
%	restrictions $\cI^{(i)}_\nu\vert_{\mathscr{U}_N}$ vanish for $\nu\notin\Lambda_N$ and are well defined polynomials in a finite number of variables for$\nu\in\Lambda_N$.\\
%	In this setting, identity \eqref{Inegative} restricted to
%	$\mathscr{U}_N$ reads
%	\[
%	{\I_{-h}}_{\big\vert_{\mathscr{U}_N}}=\sum_{\nu\in\Lambda_N}
%	{\Big(-\nu^{-h-1}\cI^{(1)}_\nu+(h+1)\nu^{-h}\cI^{(2)}_\nu\Big)}_{
%	\Big\vert_{\mathscr{U}_N}},
%	\]
%	 and is no longer a purely formal identity.\\
%    Now, denote by $\mathcal{P}_{h_1,h_2,N}$ the restriction of $\mathcal{P}_{h_1,h_2}$ to
%	$\mathscr{U}_N$. As in \ref{prop:comm_gamma2}, $\mathcal{P}_{h_1,h_2,N}$ coincides with the Poisson bracket of the restrictions of $\I_{-h_1}$ and $\I_{-h_2}$.
%	By the commutation of the
%	${\cI^{(i)}_\nu}_{\vert_{\mathscr{U}_N}}$ proved in \ref{prop:comm_gamma2},
%	we have
%	\[
%	\mathcal{P}_{h_1,h_2,N}=0\,,\qquad \forall N\in\N\,.
%	\]
%	Since $\mathcal{P}_{h_1,h_2}$ is continuous on $\H^1$ and
%	$\lim_{N\to\infty}\mathcal{P}_{h_1,h_2,N}=\mathcal{P}_{h_1,h_2}$, we conclude.
\end{proof}
%Recalling
%\[
%\frac{x}{x-1} = -\sum_{n=1}^{\infty} x^n\,, \quad \frac{x}{(x-1)^2} = x \cdot \sum_{n=1}^{\infty} n x^{n-1} = \sum_{n=1}^{\infty} n x^n
%\]
%we have
%\[
%\mathcal J(\frac1\kappa)= \sum_{\nu \in \Lambda} \sum_{n=1}^{\infty}  -(\kappa \nu^{-1}))^n \cI_\nu^{(1)} + n \nu (\kappa \nu^{-1})^n \cI_\nu^{(2)} 
%\]
%\[
%\mathcal J(\frac1\kappa)=  \sum_{h=0}^{\infty} \kappa^{h+1} \sum_{\nu\in \Lambda} \Big(- \nu^{-h-1}\cI_\nu^{(1)} + n \nu^{-h} \cI^{(2)}_\nu\Big) 
%\]
\section{The Kirchhoff-Pohozaev hierarchy in Fourier/complex notation}\label{complessi}


%\red{l'equazione qui sopra va spostata nell'introduzione: non lo faccio adesso per non creare casino}

%Then system \eqref{sysJ} becomes
%\begin{equation}\label{u v fourier}
%\begin{cases}
%\partial_t u_j = v_j \\
%\partial_t v_j = - \dfrac{|j|^2 u_j}{\left(1 + \sum_{k\in\Z^n\setminus\{0\}} |k|^2 |u_k|^2 \right)^2},
%\end{cases}
%\qquad j\in\Z^n\setminus\{0\}.
%\end{equation}

In this section, we introduce the standard complex coordinates that diagonalize the quadratic part of the Hamiltonian, and we study the expression of the constants of motion in terms of the complex variables.
%As usual, we symmetrize the order of derivation and diagonalize the system linearized at zero by introducing the complex coordinates
Recalling \eqref{fourier}, we define
\begin{equation}\label{z bar z}
z_j := \frac{|j|^{\frac12} u_j + i |j|^{-\frac12} v_j}{\sqrt2}, \qquad
\bar z_j := \frac{|j|^{\frac12} u_{-j} - i |j|^{-\frac12} v_{-j}}{\sqrt2}.
\end{equation}
%which is equivalent to
%\begin{equation}\label{z inv}
%u_j = \frac{|j|^{-\frac12} (z_j + \bar z_{-j}) }{\sqrt2}, \qquad v_j = \frac{|j|^{\frac12} (z_j - \bar z_{-j}) }{i\sqrt2}.
%\end{equation}
%In the complex variables, system \eqref{u v fourier} becomes
%\begin{equation}\label{system z}
%\begin{cases}
%\partial_t z_j = -i|j|z_j -\dfrac{i(z_j + \bar z_{-j})}{2\left( 1 + \frac12\sum_{k\in\Z^n\setminus\{0\}} |k|(z_k+\bar z_{-k})(\bar z_k+z_{-k}) \right)^2}\\[1cm]
%\partial_t \bar z_j = i|j| \bar z_j + \dfrac{i(\bar z_j + z_{-j})}{2\left( 1 + \frac12\sum_{k\in\Z^n\setminus\{0\}} |k|(z_k+\bar z_{-k})(\bar z_k+z_{-k}) \right)^2}
%\end{cases}
%\end{equation}
%In these variables, the Hamiltonian \eqref{hamiltonian1int} reads
%\begin{equation}\label{ham z}
%\begin{aligned}
%	H(z,\bar z) =& -\frac12 + \sum_{j\in\Z^n\setminus\{0\}} |j| |z_j|^2
%- \dfrac{\left( \sum_{j\in\Z^n\setminus\{0\}} |j| (|z_j|^2 + |z_{-j}|^ 2 + z_j z_{-j} + \bar z_j \bar z_{-j}) \right)^2}{8+ 4 \sum_{j\in\Z^n\setminus\{0\}} |j| (|z_j|^2 + |z_{-j}|^ 2 + z_j z_{-j} + \bar z_j \bar z_{-j}) }
%\\
%=& -\dfrac12 +\sum_{j\in\Z^n\setminus\{0\}} |j| |z_j|^2
%- \frac12 \dfrac{\left( \sum_{j\in\Z^n\setminus\{0\}} |j| (|z_j|^2 + {\rm Re}z_j z_{-j} ) \right)^2}{1+ \sum_{j\in\Z^n\setminus\{0\}} |j| (|z_j|^2 + {\rm Re} z_j z_{-j}) }.
%\end{aligned}
%\end{equation}
%
%\[\begin{aligned}
%	D_k &= \sum_{j \neq 0} |j|^{2k-1} \left( |z_j|^2 + \text{Re}(z_j z_{-j}) \right) \,,\quad 
%	B_k = \sum_{j \neq 0} |j|^{2k-1} \left( |z_j|^2 - \text{Re}(z_j z_{-j}) \right)  \\
%	E_k &= \sum_{j \neq 0} |j|^{2k-1} \left[ \left(\frac{1}{A} + A\right)|z_j|^2 - \left(\frac{1}{A} - A\right)\text{Re}(z_j z_{-j}) \right] \,,\quad 
%	Q_k = \sum_{j \neq 0} |j|^{2k-2} \text{Im}(z_j z_{-j}) \\
%	F_k &= \sum_{j \neq 0} |j|^{2k-1} \left[ \left(\frac{1}{A} - A\right)|z_j|^2 - \left(\frac{1}{A} + A\right)\text{Re}(z_j z_{-j}) \right]
%\end{aligned}\]
It is also convenient to define
\[	\begin{aligned}
	H^{(2)}_k:= \sum_{j \neq 0} |j|^{2k-1} |z_j|^2\,,\quad P_k:= \sum_{j \neq 0} |j|^{2k-1}
	\text{Re}(z_j z_{-j})\,,
	%\quad Q_k = \sum_{j \neq 0} |j|^{2k-2} \text{Im}(z_j z_{-j})
\end{aligned}
\]
so that, in the variables $(z,\bar z)$, the quantities defined in \eqref{BDQ} read
\[
\begin{aligned}
	D_k &= H^{(2)}_k + P_k\,, \quad  B_k = H^{(2)}_k - P_k\,, \quad Q_k = \sum_{j \neq 0} |j|^{2k-2} \text{Im}(z_j z_{-j}), \\
	E_k &= 2 H^{(2)}_k + \frac{D_1^2 + 2D_1}{1+D_1} B_k - \frac{2D_1}{1+D_1} H^{(2)}_k, \\
	F_k &= -2 P_k + \frac{D_1^2 + 2D_1}{1+D_1} B_k + \frac{2D_1}{1+D_1} P_k\,.
\end{aligned}
\]
Substituting in the expressions for $H$ and $\I_k$, we obtain %that the constants of motion %(questa scrittura ora  Simone lo sposta alla sezione 1)
%\begin{equation}\label{nuova}
%	\I_k=E_k+\dfrac{1}{4}\sum_{j=1}^{k-2}(E_{1+j}E_{k-j}-F_{j+1}F_{k-j})-\sum_{j=2}^k Q_{j}Q_{k-j+2}.
%\end{equation}
% are now written as
%are
%(for $k\geq2$, we divide them by two for convenience)
\begin{equation}\label{costanticompl}
	\begin{aligned}
		H &= -\frac12 + H_1^{(2)} - \frac{D_1^2}{2(1+D_1)} \quad\text{or equivalently}\\
		\I_1&= -2H_1^{(2)} + \frac{D_1^2}{1+D_1},\\[15pt]
		\I_k &= 2H^{(2)}_k + \frac{D_1^2 + 2D_1}{1+D_1} B_k - \frac{2D_1}{1+D_1} H^{(2)}_k \\%[10pt]
		&\quad + \sum_{j=1}^{k-2} \left( H^{(2)}_{1+j}H^{(2)}_{k-j} - P_{1+j}P_{k-j} \right) - \sum_{j=2}^k Q_{j}Q_{k-j+2} \quad\  \forall k\geq2.
	\end{aligned}
\end{equation}
Focusing on the quadratic part of the constants of motion, we have the following lemma.
\begin{lemma}\label{lemma complex const}
	One has
	\begin{equation}\label{quadratic constants z}
		\begin{aligned}
			\I_1&= -2H_1^{(2)} + \I_1^{(\geq4)}, \qquad \I_k= 2H_k^{(2)} + \I_k^{(\geq4)} \quad\  \forall k\geq2,\\
			\cI_\nu^{(1)} &= 2\sqrt{\nu} \sum_{j\in S_\nu} |z_j|^2 + \cI_\nu^{(1, \geq4)} \quad\ \forall\nu\in\Lambda,
		\end{aligned}
	\end{equation}
	where $\I_1^{(\geq4)}, \I_k^{(\geq4)}, \cI_\nu^{(1, \geq4)}$ are of polynomial homogeneity $\geq4$ as functions of the variables $(z,\bar z)$.
\end{lemma}
\begin{proof}
	For $\I_1, \I_k$, the thesis follows directly from \eqref{costanticompl}. For $\cI_\nu^{(1)}$, from the definition in \eqref{eq:Gamma1} it is enough to note that $A=(1+D_1)^{-1}$ and therefore the quadratic part of $\cI_\nu^{(1)}$ is just $\sum_{j \in S_\nu} (\nu|u_j|^2 + |v_j|^2)$. The thesis follows by a straightforward computation by applying the change of coordinates \eqref{z bar z}.
\end{proof}








%To write the generating function we  define
%\[
%\begin{aligned}
%	&H^{(2)}(\lambda):= \sum_{k=2}^\infty \lambda^{k-1} \sum_{j \neq 0} |j|^{2k-1} |z_j|^2\,,\quad P_k:= \sum_{k=2}^\infty  \lambda^{k-1} \sum_{j \neq 0} |j|^{2k-1}
%	\text{Re}(z_j z_{-j})\\ &
%	Q_k = \sum_{k=2}^\infty \lambda^{k-1}  \sum_{j \neq 0} |j|^{2k-2} \text{Im}(z_j z_{-j})
%\end{aligned}
%\]
%
%recalling that
%\[
%\sum_{k=2}^\infty \lambda^{k-1} |j|^{2k-1} = \frac{\lambda |j|^3}{1 - \lambda |j|^2}
%\]
%
%we obtain at least formally
%
%\[
%\begin{aligned}
%	H^{(2)}(\lambda) &= \sum_{j \neq 0} \frac{\lambda |j|^3}{1 - \lambda |j|^2} |z_j|^2 \,,\quad 
%	P(\lambda) &= \sum_{j \neq 0} \frac{\lambda |j|^3}{1 - \lambda |j|^2} \text{Re}(z_j z_{-j}) \\[10pt]
%	Q(\lambda) &= \sum_{j \neq 0} \frac{\lambda |j|^2}{1 - \lambda |j|^2} \text{Im}(z_j z_{-j})
%\end{aligned}
%\]
%then substituting in $\I(\lambda)$ we obtain the expression:
%\[
%\begin{aligned}
%	\I(\lambda) &= \frac{D_1^2 + 2D_1 + 2}{1+D_1} H^{(2)}(\lambda) - \frac{D_1^2 + 2D_1}{1+D_1} P(\lambda) \\[10pt]
%	&\quad - \frac{1}{\lambda}Q^2(\lambda) + \left(H^{(2)}(\lambda)\right)^2 - \left(P(\lambda)\right)^2
%\end{aligned}
%\]
%which reads:
%\[
%\begin{aligned}
%	\I(\lambda) &= \left(2 + \frac{D_1^2}{1+D_1}\right)  \sum_{j \neq 0} \frac{\lambda |j|^3}{1 - \lambda |j|^2}  |z_j|^2 - \frac{D_1^2 + 2D_1}{1+D_1}  \sum_{j \neq 0} \frac{\lambda |j|^3}{1 - \lambda |j|^2}\text{Re}(z_j z_{-j}) \\[15pt]
%	&\quad + \sum_{j \neq 0} \frac{\lambda^2 |j|^6}{(1 - \lambda |j|^2)^2} \left( |z_j|^4 - |z_j z_{-j}|^2 \right)
%\end{aligned}
%\]
%recalling that the equation has $M_j:= |z_j|^2-|z_{-j}|^2$  as constants of motion we get
%\[
%\begin{aligned}
%	\I(\lambda) &= \left(2 + \frac{D_1^2}{1+D_1}\right)  \sum_{j \neq 0} \frac{\lambda |j|^3}{1 - \lambda |j|^2}  |z_j|^2 - \frac{D_1^2 + 2D_1}{1+D_1}  \sum_{j \neq 0} \frac{\lambda |j|^3}{1 - \lambda |j|^2}\text{Re}(z_j z_{-j}) \\[15pt]
%	&\quad + \frac{1}{2} \sum_{j \neq 0} \frac{\lambda^2 |j|^6}{(1 - \lambda |j|^2)^2} M_j^2= 
%	E(\lambda)+ \frac{1}{2} \sum_{j \neq 0} \frac{\lambda^2 |j|^6}{(1 - \lambda |j|^2)^2} M_j^2
%\end{aligned}
%\]
%


\section{On Integrability of the Kirchhoff-Pohozaev hierarchy.}
\subsection{Restriction to finite dimensional support}
As explained in the introduction, due to the constants of motion $M_j$, for any symmetric $T\subset \Z^n\setminus\{0\}$ the sets
\[
\mathscr{U}_T:=\{(u,v)\in \H^1 \;\vert{} \quad u_j=v_j=0 \,,\quad \forall j\notin T \}
\]
defined in \eqref{UT} are invariant under the dynamics. Let us assume that $T$ has finite cardinality and write w.l.o.g, $T=\cup_{h=1}^N T_h$ where $T_h= T\cap S_{\nu_h}\neq \emptyset$ for some  increasing sequence $\nu_h\in \Lambda$ and consider the functions $\cI,M$ restricted to $\mathscr{U}_T$.
\begin{lemma}\label{restringo}
	The number $A$ of of algebraically independent  resatricted constants of motion $\cI,M$ is 
	\[
	N+|T|\leq A\leq N + |T|/2 + K 
	\]
	where $K$ is the number of $T_h$ with cardinality $\geq 4$. In particular in dimension $n=1$ the number of independent constants of motion is equal to $|T|$ and the restricted system is integrable. 
\end{lemma}
\begin{proof}
	This is a direct consequence of the definitions. We start by noticing that all the restricted $M_j$ with $j\notin T$  are zero as well as all the restricted $\cI^{(a)}_{\nu}$ such that $T\cap S_{\nu}=0$. Recalling that $M_j=M_{-j}$ we divide 
	$T= T_+\cup T_{-}$ with $T_{-}=-T_{+}$ disjoint.  We 
	consider  the list $\{M_j\}_{j\in T_+}$, of course they are independent since the have different variables. Now we add the list $\{\cI^{(1)}_{\nu_h}\}_{h=1}^N$, these are all non-trivial polynomials. Finally we remark that
	\[
	\cI_{\nu}^{(2)}= 2\sum_{j\in S_\nu} M_j^2 + \frac12 \sum_{\substack{j,k\in S_\nu\\ j\neq \pm k}}|u_j v_k -v_j u_k|^2
	\]
	thus, if $\{h_i\}_{i=1}^K$ is the increasing sequence  of indexes for which $|T_{h_i}|\geq 4$, we define $$L_i:= \sum_{\substack{j,k\in S_{\nu_{h_i}}\\ j\neq \pm k}}|u_j v_k -v_j u_k|^2$$
	In conclusion our candidates for independent constants of motion are the restrictions to $\mathscr{U}_T$ of
	\begin{equation}
		\label{costantine}
		\{M_j\}_{j\in T_+}\cup\{\cI^{(1)}_{\nu_h}\}_{h=1}^N\cup  \{ L_i \}_{i=1}^K
	\end{equation}
	the fact that  the list $\{M_j\}_{j\in T_+}\cup\{\cI^{(1)}_{\nu_h}\}_{h=1}^N$ is made of algebraically independent functions follows since they are polynomials and their quadratic parts are (written in the complex variables $z_j$ of section \ref{complessi})
	\[
	M_j=(M_j)^{\rm quadratic} = |z_j|^2-|z_{-j}|^2\,,\quad (\cI_{\nu_h}^{(1)})^{\rm quadratic} = \sqrt{\nu_h}\sum_{j\in T_h} |z_j|^2
	\]
	which are evidently independent. Regarding the $L_i$, it is obvious that the are independent from each other and from the $M_j$, however it is not so trivial to control the indepedndence from the $\cI^{(1)}$ (we strongly believe this to be true, however...)
\end{proof}
We are now ready to prove Theorem \ref{finito}.
\begin{proof}[Proof of Theorem \ref{finito}. ]
	We start by noticing that, if $$T=\{j_1,-j_1,\dots, j_d,-j_d\}\subset  \Z^n\setminus\{0\}$$ then $K=0$ and the constants of motion in \eqref{costantine} reduce to \eqref{massimali?}. Since these are $2d$ independent constants of motion in involution we have an integrable system. Moreover the Hamiltonian \ref{hamiltonian1int}, by Lemma \ref{calIJ} satisfies
	\[
	H= \frac12 (\sum_{j=1}^d \cI^{(1)}_{|j_d|} -1)\,,
	\]
	hence it is integrable. This concludes the proof of item $(i)$. In order to prove item $(ii)$ we just apply Vey's theorem (see \cite{Vey1978},\cite{Eliasson1990}) in the version proposed for example by \cite{Ito1989} or \cite{Stolovitch2000}.
\end{proof}
\subsection{ Formal Birkhoff  Normal Form on sequence spaces}
In the paper \cite{ProcesiStolo} there is a discussion of the Poisson algebra of formal Hamiltonians with a fixed point at zero, on formal symplectic changes of variables which preserve the fixed point and on the existence of a Formal Birkhoff normal form. Our purpose here is to show that formal Hamiltonians in involution can be put in simultaneous normal form, and how the presence of commuting constants of motion simplfies the normal form itself. This is a classical argument in finite dimension, see for example \cite{Ito1989}, we briwfly review the argument as well as those resulta  nneded to apply it in infinte dimension.
\\
Even though we could work in the real variables, it is convenient to pass to the complex notations of the previous section. Here, for the first time, we shall strongly need to restrict to the one dimensional torus, hence set  $n=1$ so that $j\in \Z\setminus\{0\}$.
Let $I\subset \Z$ be an index set (we shall be only interested in the case $I=\Z_0=\Z\setminus\{0\}$).
As usual given a vector $k\in \Z^I$, $|k|:=\sum_{j\in I}|k_j|$. We denote $\N^I_f$ to be the set of  finitely supported sequences of non negative integers, similarly for $\Z^I_f$. If $j\in\mathbb{I}$ then $\be_j\in \Z^I_f$ denotes the vector the $j$-coordinate of which is  $1$, while the others are zero.
%\com{The definition of formal power series does not depend in any way on $s,p,\theta$ but only on the fact that $u=\pa{z_j}_{j\in\Z} \in \C^\Z$}
%I am systematically assuming that $\N$ contains zero.

\begin{defn}[Formal power series]\label{Hr}
	We consider the space $\cF$ of formal  power series expansions  %\footnote{As usual given a vector $k\in \Z^\Z$, 		$|k|:=\sum_{j\in\Z}|k_j|$.} 
	in $z=(z_j)_{j\in I}\in \C^I$:
	$$ 
	H(u)  = \sum_{\substack{\al,\bt\in\N^I_f} }H_{\al,\bt}z^\al \bar z^\bt\,,
	\qquad u\in \C^\Z,\quad
	z^\al:=\prod_{j\in I}z_j^{\al_j} 
	$$ 
	with the following properties: 
	\begin{enumerate}%[(i)]
		\item $H_{0,0}= 0$, $H_{\be_0,0}= H_{0,\be_0}=0$ %\,,\quad\xout{ $H_{\be_i,0} =0$ for all $\be_i$ }\com{ this last condition follows by momentum conservation}
		\item Reality condition:
		\begin{equation}\label{real}
			H_{\al,\bt}= \overline{ H}_{\bt,\al}\,;
		\end{equation}
		\item Momentum conservation:
		\begin{equation}
			H_{\al,\bt}= 0 \quad\mbox{if}\;\, \pi(\al,\bt):= \sum_{j\in I} j(\al_j-\bt_j)\ne 0 \label{mconserv}
		\end{equation}%namely the Hamiltonian Poisson commutes with the { \sl momentum} $\sum_{j\in \Z}j |z_j|^2$;
	\end{enumerate}
	We shall denote 
	\[
	\cM:=\{(\al,\bt) \in \N^I_f:\qquad \pi(\al,\bt)=0\}
	\]
	so that $H\in \cF$ can be written as
	\[
	\sum_{(\al,\bt)\in \cM}H_{\al,\bt}z^\al \bar z^\bt
	\] 
	
\end{defn}

\begin{lemma}\label{anello}
	The space $\cF$ is a ring  with respect to the  sum and multiplication
	\[
	\begin{aligned}
		(H+F)(z) &:=  \sum_{(\al,\bt)\in \cM} (H_{\al,\bt}+ F_{\al,\bt}) z^\al \bar z^\bt \\
		(H \cdot F )(z) &:=  \sum_{(\al,\bt)\in \cM} \Big(\sum_{(\al_1,\bt_1) \in \cM} \sum_{\substack {(\al_2,\bt_2) \in \cM \\
				\al_1+\al_2=\al\,,\\ \bt_1+\bt_2=\bt}} H_{\al_1,\bt_1} F_{\al_2,\bt_2}\Big) z^\al \bar z^\bt\,.
	\end{aligned}
	\]
	Moreover $\cF$ is a Poisson albebra with respect to the Brackets
	\begin{equation}\label{well}
		\set{F,G} := \im\sum_{  (\al^{(i)},\bt^{(i))}\in\cM }  \bcoeffu{F}\bcoeffd{G} \sum_j \pa{\aluno_j \btdue_j  - \btuno_j \aldue_j }z^{\aluno + \aldue -\be_j}\bar{u}^{\btuno + \btdue - \be_j}\,.
	\end{equation}
	
\end{lemma}
\begin{proof}
	We only need to show that the definitions are  well posed, which is  obvious for the sum and, for the multiplication, is due to the fact that, since $\al_i\bt_i\geq 0$ the sums over these indexes are finite. Regarding the Poisson brackets, this is proved in \cite{ProcesiStolo}.
\end{proof}

Let us set
\begin{equation}\label{nocciolina}
	\cK:=\left\{Z\in \cF\, : \, Z(z) = \!\!\sum_{\substack{\al\in \cM:}} Z_{\al,\al}|u|^{2\al}\right\}\,,\qquad \cR:=\left\{R\in \cF\, : \, R(z) = \!\!\!\!\!\!\!\!\sum_{\al,\bt\in \cM:\, \al\ne \bt }\!\!\!\!\!\!\!\! R_{\al,\bt}z^\al\bar z^\bt\right\}
\end{equation}


%\com{what does "momentum conservation" mean in finite dimension  ? it seems to be very restrictive ! We need to show that nonlineartity cannot have $\partial_x$ terms, that is of the form $\dot u= F(u)$}
%\blu{In the what follow and unless it is necessary, we shall omit to indicate the indices $s,p,\theta$.}
\begin{defn}[ scaling degree]\label{scialla}
	For $d\in \N$, we denote by $\cF^\td\subset \cF$ the vector space of homogeneous formal polynomials of degree $\td+2$, and define 
	%\[
	%	\xout{	\cF^{\leq d}= \cup_{h\leq d}\,,\quad \cF^{>d}:=\cF\setminus \cF^{\leq d} \,,\quad  \cF^{\ge d}:= \cF^{>d}\cup \cF^{d}\,,\dots}
	%	\]
	\[
	\cF^{\leq d}= \oplus_{h\leq d} \cF^{h}\,,\quad \cF^{>d}:=\widehat\oplus_{h>d} \cF^{h} \,,\quad  \cF^{\geq d}:= \cF^{>d}\oplus \cF^{d}\,,\cF =\cF^{\leq d}\oplus\cF^{> d},\dots
	\]
	We define the projections  associated to these direct sum decompositions
	\[
	\Pi^{(\td)} H= \sum_{|\al|+|\bt|=\td+2}H_{\al,\bt} z^\al\bar z^\bt \,,\quad \Pi^{(> \td)} H= \sum_{|\al|+|\bt|> \td+2}H_{\al,\bt} z^\al\bar z^\bt\,,\dots
	\]
	Elements of $\cF^{\geq d}$ (resp. $\cF^{> d}$) are said to be {\it of scaling degree} $\geq d+2$ (resp. $>d+2$).
	Finally we define
	\[
	\Pi^\cK H =  \sum_{\al}H_{\al,\al} |z|^{2\al}\,,\quad \Pi^\cR H = \sum_{\al\ne \bt}H_{\al,\bt} z^\al\bar z^\bt.
	\]
	We denote by $\cK^d:= \cF^d \cap \cK$  and similarly for $\cR$ and $\geq \td ,\leq \td$.
	Note that $\cF= \widehat \oplus_d \cF^d$.
\end{defn}

\begin{rmk}
	By construction  the scaling degree is well behaved under multiplication and Poisson brackets. Indeed, if one has $H\in\cF^{\td_1}$, $K\in\cF^{\td_2}$  then $H K\in\cF^{\td_1+\td_2+2}$ and $\{H,K\} \in\cF^{\td_1+\td_2}$.
	%	Of course, since we are in infinite dimension even if  the $\cF^d$ are homogeneous,  they are  only formal polynomials.  However, if we restrict to monomials $z^\al\bar{u}^\bt$ with $|\al_j|+|\bt_j| =0$ for all $|j|>N$,  we are working on the usual space of polynomials, on which we have the standard symplectic structure $\im \sum_{j\leq N} d z_j\wedge d \bar z_j$. We now  show that such structure extends to $\cF$.
\end{rmk}
%\begin{prop}\label{degree decompositionbis}
%	The following Formula \eqref{well} is well defined and  endows $\cF$ with a  Poisson algebra structure  which is a filtered  Lie algebra w.r.t. the $\cF^{\geq d}$'s. 
%	\begin{equation}\label{well}
	%		\set{F,G} := \im\sum_{  (\al^{(i)},\bt^{(i))}\in\cM }  \bcoeffu{F}\bcoeffd{G} \sum_j \pa{\aluno_j \btdue_j  - \btuno_j \aldue_j }z^{\aluno + \aldue -\be_j}\bar{u}^{\btuno + \btdue - \be_j}
	%	\end{equation}
%\end{prop}
%Before proving our assertion we need a technical lemma.
%Let $\be_j\in \N^I_f$ be the $j$th vector of the standard basis.
%\begin{lemma}\label{finmo} 1)\quad  Given $\alpha\in \N^I_f$ there  is only a finite number of pairs    $ \alpha^{(1)},\alpha^{(2)}\in \N^I_f$ with $\alpha=\alpha^{(1)}+\alpha^{(2)}$.
%	2)\quad 	
%	Given  $ (\al,\bt)\in \cM$ there is only a finite number of pairs    $ (\aluno,\btuno),\  (\aldue,\btdue)\in \cM$ and indices $j\in\Z$  such that:
%	\begin{enumerate} 
	%		\item[\it i)] $ (\al,\bt)= (\aluno,\btuno)+(\aldue,\btdue)-(\be_j,\be_j)$ 
	%		\item[\it ii)] one has $\al_j^{(1)}\bt_j^{(2)}+ \al_j^{(2)}\bt_j^{(1)}\ne 0$.%the index $j$ satisfies at least one of the two following conditions,    $(\aluno=\bar \aluno+\be_j ,\btdue=\bar \btdue+\be_j)$ or $ (\btuno=\bar \btuno+\be_j,\aldue= \bar \aldue+\be_j )$ with $\bar\al_1,\bar{\bt}_1\in \N^I_f$.
	%	\end{enumerate}  
%\end{lemma} 
%\begin{proof} 1) is clear since for all $j$ one has $0\leq (\alpha_1)_j\leq  \alpha_j  $.\smallskip
%	\\
%	2)\quad By item 1) we may divide $(\al,\bt)= (a^{(1)}, b^{(1)})+ (a^{(2)}, b^{(2)})$ in a finite number of ways. Then 
%	the pairs 	$(\aluno,\btuno),(\aldue,\btdue)$  can only have one of the following forms (up to exchanging the indices)
%	
%	A)	$\quad (\aluno,\btuno)= (a^{(1)}, b^{(1)})+ (\be_j,\be_j)\,,\quad (\aldue,\btdue) = (a^{(2)}, b^{(2)})$
%	
%	B) $\quad (\aluno,\btuno)= (a^{(1)}, b^{(1)})+ (\be_j,0)\,,\quad (\aldue,\btdue) = (a^{(2)}, b^{(2)})+(0,\be_j)$,
%	\\
%	for some index $j\in \Z$.
%	\\
%	If we are in case $A)$ then 	by condition $ii)$  we have $j\in $Supp$(a^{(2)}+b^{(2)})$, which restricts to a finite number of possible $j's$. Otherwise in case $B)$ 
%	by momentum conservation e have $j= -\pi( a^{(1)}, b^{(1)})={\pi} (a^{(2)}, b^{(2)})$ and again $j$ is restricted to a finite number of possible choices.
%\end{proof}
%\begin{proof}[Proof of Proposition \ref{degree decompositionbis}]
%	The fact that the Poisson bracket is well defined  follows immediately from the previous Lemma. 
%	Indeed by construction
%	\[
%	\{F,G\}= \sum_{\al,\bt} P_{\al,\bt} z^\al\bar{z}^\bt\in \cF
%	\]
%	where  $P_{\al,\bt}=0$ if $\pi(\al,\bt)\ne 0$ and otherwise
%	\begin{equation}
	%		\label{begon}
	%		P_{\al,\bt}= {\im}\sum_j \sum_{\substack {\al^{(i)},\bt^{(i)}\in\N^I_f:\; \pi(\al^{(i)},\bt^{(i)})=0\\ \al = \aluno+\aldue - \be_j\,, \bt = \btuno+\btdue - \be_j }}
	%		\bcoeffu{F}\bcoeffd{G} \pa{\aluno_j \btdue_j  - \btuno_j \aldue_j }\,.
	%	\end{equation}
%	Then item 2 of the previous Lemma  implies that $P_{\al,\bt}$ above is given by a finite sum.
%	\\
%	The fact that it  endows $\cF$ with a  Poisson algebra structure follows from the fact that   the infinitely many identities defining such a structure   involve  only a finite number of elements $z_i,\bar z_i$  and then we are in the canonical  Poisson algebra.
%	\\
%	The filtered Lie algebra property comes from the fact that in \eqref{begon} we get $|\alpha|+|\beta| = |\aluno|+ |\aldue|+ |\btuno|+ |\btdue| -2$, this shows that if $F\in \cF^{\geq \td_1}$, and $G\in \cF^{\geq \td_2}$  then 
%	\begin{equation}
	%		\label{a+b}|\alpha|+|\beta| \geq \td_1+2+\td_2+2-2 = \td_1+\td_2+2.
	%	\end{equation}  so 
%	$\{F,G\}\in \cF^{\geq \td_1+\td_2}$.
%\end{proof}




%\begin{prop}\label{degree decomposition}
%	The  Poisson Bracket defined by $\Omega$:
%	\begin{equation}
	%		\set{F,G} := \im\sum_{\substack {\al^{(i)},\bt^{(i)}\in\N^I\,, \\
			%				|\al^{(i)}|+|\bt^{(i)}|<\infty\,,\, \pi(\al^{(i)},\bt^{(i)})=0}}  \bcoeffu{F}\bcoeffd{G} \sum_j \pa{\aluno_j \btdue_j  - \btuno_j \aldue_j }z^{\aluno + \aldue -\be_j}\bar{u}^{\btuno + \btdue - \be_j}
	%	\end{equation}
%	endows $\cF$ with a graded Poisson algebra structure w.r.t. the $\cF^d$. 
%	
%\end{prop}
%\begin{proof}
%	In degree to show  that the Poisson brackets are well defined at the level of formal power series, we have to show that 
%	\[
%	\{F,G\}= \sum_{\al,\bt} P_{\al,\bt} \buu \in \cF
%	\]
%	A direct computation shows that $P_{\al,\bt}=0$ if $\pi(\al,\bt)\ne 0$ and otherwise
%	\begin{equation}
	%		\label{begon}
	%		P_{\al,\bt}= \sum_j \sum_{\substack {\al^{(i)},\bt^{(i)}\in\N^I\,, \\
			%				|\al^{(i)}|+|\bt^{(i)}|<\infty\,,\, \pi(\al^{(i)},\bt^{(i)})=0\\ \al = \aluno+\aldue - e_j\,, \bt = \btuno+\btdue - e_j }}
	%		\bcoeffu{F}\bcoeffd{G} \pa{\aluno_j \btdue_j  - \btuno_j \aldue_j }\,.
	%	\end{equation}
%	
%	Note that $|\alpha|+|\beta| = |\aluno|+ |\aldue|+ |\btuno|+ |\btdue| -2$, this shows that if $F\in \cF^{\td_1}$, and $G\in \cF^{\td_2}$  then 
%	\begin{equation}
	%		\label{a+b}|\alpha|+|\beta| = \td_1+2+\td_2+2-2 = \td_1+\td_2+2.
	%	\end{equation} 
%	In conclusion if the sum in \eqref{begon} is well defined then 
%	$P\in \cF^{\td_1+\td_2}$.
%	\\
%	In degree to show that $P_{\al,\bt}$ is always well defined we need to show that the sums of the left hand side are finite sums.
%	\\
%	Let $\alpha,\beta\in \N^\Z$ such that $|\alpha|,|\beta|<+\infty$ and $\pi(\alpha,\beta)=0$. 
%	Recalling that since $|\al|<\infty$,  there are a finite number of ways of dividing $\al= a^{(1)}+ a^{(2)}$ with the $a^{(i)}\in \N^\Z$. The same holds for $\bt$.  Given a quartuple $q=\{a^{(1)}, a^{(2)},b^{(1)}, b^{(2)}\}$ of elements of $\N^\Z$ with $\al= a^{(1)}+ a^{(2)}$ and
%	$\bt= b^{(1)}+ b^{(2)}$ and recalling that $\pi(\al-\bt)= 0$, we have 
%	$$
%	0=\sum_ii\left((a^{(1)}_i+a^{(2)}_i)-(b^{(1)}_i+b^{(2)}_i)\right).
%	$$
%	Hence,	we have 
%	\[
%	j(q):=	\sum_i i(a_i^{(2)}- b_i^{(2)})= -\sum_i i(a_i^{(1)}- b_i^{(1)}).
%	\]
%	We define
%	$$
%	\al^{(1)}(q)= a^{(1)} + \be_{j(q) }\,,\;  \bt^{(1)}(q)= b^{(1)},\quad \al^{(2)}(q)= a^{(2)}\,,\;  \bt^{(2)}(q)= b^{(2)} + \be_{j(q) },
%	$$
%	Hence, we have $\pi(\al^{(i)},\bt^{(i)})=0$, $i=1,2$,  $\alpha=\al^{(1)}(q)+\al^{(2)}(q)- \be_{j(q) }$ and $\bt=\bt^{(1)}(q)+\bt^{(2)}(q)- \be_{j(q) }$. We remark also that if we set $\tilde q=\{a^{(2)}, a^{(1)},b^{(2)}, b^{(1)}\}$, then $j(\tilde q)=-j(q)$.
%	%On the other hand, for any $c^{(i)}\in \N^{\Z}$, we can change $\al^{(i)}(q),\bt^{(i)}(q)$ into $\al^{(i)}(q)+c^{(i)},\bt^{(i)}(q)+c^{(i)}$
%	Hence, $P_{\al,\beta}$ is defined as it is expressed as a finite sum~: 
%	$$
%	P_{\al,\bt}=\sum_{q=\{a^{(1)}, a^{(2)},b^{(1)}, b^{(2)}\}}F_{a^{(1)} + \be_{j(q) }, b^{(1)}}G_{a^{(2)}, b^{(2)}+ \be_{j(q) }}\pa{(a^{(1)}_{j(q)}+1)(b^{(2)}_{j(q)}+1)  - a^{(2)}_{j(q)}b^{(1)}_{j(q)} }
%	$$
%\end{proof}
\begin{rmk}\label{gogna}
	Let $H_i\in \cF^{\geq \td_i}$ be a sequence of formal Hamiltonians with  $\td_{i+1}> \td_i$ for all $i\geq 1$. Then the series
	\[
	H= \sum_{i=1}^\infty H_i \in \cF^{\geq \td_1}
	\]
	is well defined since for any $\td\geq \td_0$ the projection
	\[\Pi^{(\leq \td)} H = \Pi^{(\leq \td)} \sum_{i: \td_i \leq \td}  H_i  \] is a finite sum.
\end{rmk}
We say that  a linear operator $L:\cF\to \cF$ is of degree (or increase the degree by) $\td$  if for all $h$ 
\[
L:\cF^{ \geq h} \to \cF^{ \geq h+\td}\,.
\]
We now state some Lemmas, regarding formal power series in infinitely many variables, which are proved in \cite{ProcesiStolo} and are used to ensure that a formal Birckhoff Normal Form exists.
\begin{lemma}\label{lemniscata}
	let $L_n$ be a sequence of linear operators on $\cF$  and let  $\td_n$ be the degree of $L_n$. If the sequence $\td_n$ increases to infinity then 
	\[
	L:= \sum_{n=1}^\infty L_n \,,\qquad T =\prod_{n=1}^\infty (\id + L_n)-\id 
	\]
	are  linear operators on $\cF$ of degree  $\td_1$.
	%\end{lemma}
	%\begin{proof}
	%	For the first statement, for all $\td\in \N$  let   $N(\td)$  be the largest $N$ such that $\td_N\leq \td$. By construction
	%	$\Pi^{(\leq \td)} L_n K=0$ for all $n> N(\td)$ and for any $K\in \cF$.
	%	Then for all $K\in \cF$ and $N>N(\td)$ one has
	%	\[
	%	\Pi^{(\leq \td)}	\sum_{n=1}^N L_n K  = \Pi^{(\leq \td)}	\sum_{n=1}^{N(\td)} L_n K\,,
	%	\] and the claim follows.
	%	\\
	%	Regarding the second statement we proceed similarly
	%	\[
	%	\prod _{n=1}^N (\id + L_n ) =  \prod _{n=1}^{N-1} (\id + L_n ) +L_N \prod _{n=1}^{N-1} (\id + L_n )\,,
	%	\] hence, {for all $\td\geq 0$ and all $N> N(\td)$}
	%	\[
	%	\Pi^{(\leq \td)}\prod _{n=1}^N (\id + L_n ) =  \Pi^{(\leq \td)}\prod _{n=1}^{N(\td)} (\id + L_n )\,.
	%	\]
	%\end{proof} As a direct consequence we have the following.
	%\begin{cor}
	\\
	Given $G\in \cF^{\geq  \td}$, with $\td\geq 1$  we define 
	\begin{equation}
		\label{flusso}
		\ad_G:= \{G,\cdot \} \,,\quad \Phi_G:=	\exp(\{G,\cdot\})= \sum_{k\geq 0} \frac{\ad_G^k}{k!}\,,
	\end{equation}
	then $ 	\ad_G$ and  $\Phi_G-\id$  are operators of degree $\td$, namely
	\[
	\ad_G,\Phi_G -\id: \cF^{\geq h} \to \cF^{\geq h+\td}\,.
	\]
	Similarly for any sequence $b_k$  one has that 
	\[
	\sum_{k\geq \tn } b_k{\ad_G^k}: \cF^{\geq h} \to \cF^{\geq h+\td\tn }\,.
	\] 
\end{lemma}
%	\rv(\[
%	P_{\al,\bt}=\sum_{q=\{a^{(1)}, a^{(2)},b^{(1)}, b^{(2)}\}} F_{\al^{(1)}(q), b^{(1)}} G_{a^2, \bt^{(2)}(q)} \al^{(1)}_j(q)\bt^{(2)}_j(q)- 
%F_{a^1, \bt^{(1)}(q)} G_{\al^{(2)}(q), b^({2})} \al^{(2)}_j(q)\bt^{(1)}_j(q)	\])
%\begin{proof}
%	By \eqref{a+b}  with $\td=\td_2$ and $h=\td_1$, the operator $\ad_G$ maps $\cF^{\ge h}\to \cF^{\ge h+\td}$ for all $h\ge 0$.
%	Iterating  $\ad_G^k$ maps $\cF^{\ge h}\to \cF^{\ge h+k \td}$ for all $h\ge 0$.	Now consider a sequence $\cL_k$ of linear operators on $\cF$, each of degree $k\td$ (i.e. mapping $\cF^{ \ge h}\to \cF^{\ge  h+k \td}$ for all $h\ge 0$ )
%	Then 
%	\[
%	\cL:\sum_{k=0}^\infty \cL_k: \cF \to \cF\,.
%	\]
%	Indeed one writes $\cF= \hat\oplus_h \cF^h$ then $\cL$  can be expressed in block form as
%	\[
%	\Pi^{(h_1)}\cL \Pi^{(h_2)}  = 	\Pi^{(h_1)}\sum_{k \leq h_1-h_2} \cL_{k}\Pi^{(h_2)}\,.
%	\]
%\end{proof}
\begin{defn}
	Given $G\in \cF^{\geq 1}$ we call the operator $\Phi_G$ defined in \eqref{flusso} a formal symplectic change of variables on $\cF$.
\end{defn}
The following Lemma, again proved in \cite{ProcesiStolo}, ensures the group structure of the formal symplectic changes of variables
\begin{lemma}[Baker-Campbell-Hausdorff]\label{BCH}
	Given   $F\in \cF^{\geq \td_1}$ and  $G\in \cF^{\geq \td_2}$, with $\td_i\geq 1$,  then there exists $K\in \cF^{\geq 1}$, such that
	\[
	e^{\{G,\}}e^{\{F,\}}=  e^{\{ K,\}} \,,\qquad  K-F-G \in  \cF^{ \geq \td_1+\td_2}
	\] 
\end{lemma}
%\begin{proof}
%	By the Baker-Campbell-Hausdorff formula (\cite{serre-Lie1}[p.29]) one has
%	\[
%	K:= %\sum_{n=1}^{\infty} F_{1,n}=\sum_{n=1}^{\infty} \frac{1}{n}\sum_{r+s=n}\left(F_{1,r,s}'+F_{1,r,s}''\right)
%	\sum_{n=1}^\infty \frac{(-1)^{n-1}}{n} \sum_{r_i+s_i>0} \frac{[G^{r_1}F^{s_1}\dots  G^{r_n}F^{s_n}]}{%\sum_j (r_j+s_j)
	%		\left(\sum_{i=1}^n(r_i+s_i)\right) \prod_{i=1}^n r_i! s_i!}
%	\]
%	where
%	%	\begin{eqnarray*}
	%		%		F_{1,r,s}'&=& \sum_{\substack{r_1+\cdots +r_m=r\\s_1+\cdots +s_{m-1}=s-1\\r_i+s_i\geq 1, r_m\geq 1}}\frac{(-1)^{m-1}}{m}\frac{\ad_G^{r_1}\ad_F^{s_1}\dots \ad_G^{r_m}F}{r_m!\prod_{i=1}^{m-1} r_i! s_i!}\\
	%		%		F_{1,r,s}''&=& \sum_{\substack{r_1+\cdots +r_{m-1}=r-1\\s_1+\cdots +s_{m-1}=s\\r_i+s_i\geq 1}}\frac{(-1)^{m-1}}{m}\frac{\ad_G^{r_1}\ad_F^{s_1}\dots \ad_F^{s_{m-1}}G}{\prod_{i=1}^{m-1} r_i! s_i!}.
	%		%	\end{eqnarray*}	
%	\begin{equation}
	%		\label{hoho}
	%		[G^{r_1}F^{s_1}\dots G^{r_n}F^{s_n}]:= \begin{cases}
		%			\ad_G^{r_1}\ad_F^{s_1}\dots \ad_G^{r_n}F  \quad \mbox{ if} \; s_n=1
		%			\\
		%			\ad_G^{r_1}\ad_F^{s_1}\dots \ad_F^{s_{n-1}} G  \quad \mbox{ if}\; s_n=0\,,\; \mbox {and} \; r_n=1
		%			\\
		%			0  \quad \quad \quad \quad  \mbox {otherwise}
		%		\end{cases}
	%		%\{\underbrace{G,\{G,\dots\{G}_{r_1},\{\underbrace{F,\{F,\dots\{F}_{s_1}, \dots \{\underbrace{G,\{G,\dots\{G}_{r_n},\{\underbrace{F,\{F,\dots\{F}_{s_n}\}\}
	%	\end{equation}
%	Recalling that   $F\in \cF^{\geq \td_1}$ and  $G\in \cF^{\geq \td_2}$, each term $\ad_G^{r_1}\ad_F^{s_1}\dots \ad_G^{r_n}F$ (resp. $\ad_G^{r_1}\ad_F^{s_1}\dots \ad_F^{s_{n-1}}G$)
%	%$[G^{r_1}F^{s_1}\dots  G^{r_n}F^{s_n}]$ 
%	is of degree $\left(\sum_{i=1}^{n}r_i\right)\td_2+(\sum_{i=1}^{n}s_i)\td_1\geq n \min(\td_1,\td_2)$. %(resp.  $\left(\sum_{i=1}^{n-1}r_i+1\right)d+\sum_{i=1}^{n-1}s_i\geq  n+\left(\sum_{i=1}^{n-1}r_i+1\right)(d-1)$). %$\sum_{i=1}^n s_i  +d \sum_{i=1}^n r_i  \ge n$
%	%	 so 
%	%	\[
%	%	F_1^{(r)}= \sum_{n=1}^r\frac{(-1)^{n-1}}{n} \sum_{\substack{r_i+s_i>0: \\\; \sum_{i=1}^n s_i  +d \sum_{i=1}^n r_i =r} }\frac{[G^{r_1}F^{s_1}\dots  G^{r_n}F^{s_n}]}{\sum_j (r_j+s_j) \prod_{i=1}^n r_i! s_i!}
%	%	\]
%	Hence setting $N(\td)$ to be the largest $N$ such that $ N \min(\td_1,\td_2) \leq \td$n 
%	\[
%	\Pi^{\leq \td} K= \Pi^{\leq \td} \sum_{n=1}^{N(\td)} \frac{(-1)^{n-1}}{n} \sum_{r_i+s_i>0} \frac{[G^{r_1}F^{s_1}\dots  G^{r_n}F^{s_n}]}{%\sum_j (r_j+s_j)
	%		\left(\sum_{i=1}^n(r_i+s_i)\right) \prod_{i=1}^n r_i! s_i!}
%	\]
%	\\
%	Moreover if $n\geq 2 $ then  the Hamiltonian in \eqref{hoho} is of degree $\geq \td_1+\td_2$, so $K-F-G \in  \cF^{ \geq \td_1+\td_2}$. 
%\end{proof}
\begin{lemma} \label{compo} Given a sequence of generating functions $G_i\in \cF^{{\geq} \td_i}$ with $\td_{i+1} > \td_i\geq 1$ then there exists $\cG\in \cF^{\geq d_1}$ such that the composition
	\[
	\prod_i	e^{\{G_i,\}} = e^{\{\cG,\}}
	\] 
\end{lemma}
\begin{proof} 
	By Lemma \ref{lemniscata} with $L_n = e^{\{G_n,\}}  -\id$ we know that  $\prod_i	e^{\{G_i,\}}$ is a well defined operator of $\cF$. 
	Using Lemma \ref{BCH} we can define $F_k\in \cF^{\geq1}$ iteratively so that
	\[
	e^{\{F_k,\cdot\}}= e^{\{G_k,\cdot\}} e^{\{F_{k-1},\cdot\}}
	\]
	since $e^{\{G_k,\cdot\}} -\id$ is of degree $\td_k$  there exists $N(\td)$ such that if $k>N(\td)$ then 
	\[
	\Pi^{(\leq  \td )}	F_k = \Pi^{(\leq  \td )}	F_{N(\td)}
	\]
	Then  $\cG= \lim_{k\to \infty} F_k$ is well defined.
\end{proof}
%	We consider 
%		\begin{equation}\label{nocciolina}
	%			\cK:=\left\{Z\in \cF\, : \, Z(u) = \!\!\sum_{\substack{\al,\bt\in \cM:\omega\cdot (\al-\bt)=0}} Z_{\al,\bt}z^{\al}\bar z^{\bt}\right\}\,,\qquad \cR:=\left\{R\in \cF\, : \, R(u) = \!\!\!\!\!\!\!\!\sum_{\al,\bt\in \cM:\, \omega\cdot (\al-\bt)\neq 0}\!\!\!\!\!\!\!\! R_{\al,\bt}z^\al\bar z^\bt\right\}
	%		\end{equation}
%		and we can decompose $\cF= \cK\oplus \cR$  as each element of $\cF$ can uniquely be expressed in term of monomials the coefficients of which is either zero or not zero. 
%		\begin{prop}\label{formal-lin}Given any formal Hamiltonian of the form $H= D_\omega + Z+ R$, 
	%			{where $Z\in \cK^{\geq 1}$} and $R\in \cF^{\geq \td}$  with $\td\geq 1$. 
	%			
	%			{\bf Formal Normal Form:} there exists $S\in  \cF^{\geq \td}$ such that
	%				\[
	%				e^{\{S,\cdot\}} H = D_\omega  + \wtZ\,,\quad \wtZ- Z \in \cK^{\geq \td}\,.
	%				\]
	%		\end{prop}
%		\begin{proof}
	%			La devo mettere?
	%		\end{proof}
%	

\subsection{Formal Birkhoff Normal Form for families of commuting Hamiltonians}	Let us  now fix $I=\Z_0=\Z\setminus\{0\}$ and let us assume that we have a family of  formal Hamiltonians of the form:
\[ 
H_k:= L_k+ R_k^{\geq 1}\in\cF \,,\qquad L_k = \sum_{j\in \Z_0} \omega_{k.j}|z_j|^2\,,\quad R_k^{\geq 1}\in \cF^{\geq 1}
\]
for $k\in \N_{\geq 1}$ ans satisfying the following


\begin{enumerate}
	\item[H1.] {\bf Involution:} For any choice of $k_1,k_2\in \N_{\geq1}$, as well as  $k\in \N_{\geq1}$ and $j\in \N_{\geq1}$ we have
	\[
	\{H_{k_1},H_{k_2}\}=0\,,\qquad \{H_k, M_j\}=0
	\]
	where  $M_j =|z_j|^2- |z_{-j}|^2$.
	\item[H2.]{\bf Independence:} 	for any $N\geq 1$  there exists $M=M(N)$ such that  $\ell\in \Z^\Z_f$ with $|\ell|_1\leq  N$  and $\ell_j=\ell_{-j}$
	satisfies 
	\begin{equation}\label{resonance}
		\omega_k\cdot \ell =0 \,,\quad \forall k=1,\dots, M(N)
	\end{equation} if and only if $\ell=0$.
\end{enumerate}
%	Let us moreover assume that the $\omega_k$ are independent, namely that
%	for any $N\geq 1$  there exists $M=M(N)$ such that  $\ell\in \Z^\Z_f$ with $|\ell|_1\leq  N$ 
%	satisfies 
%	\begin{equation}\label{resonance}
	%		\omega_k\cdot \ell =0 \,,\quad \forall k=1,\dots, M(N)
	%	\end{equation} if and only if $\ell=0$.
%	and any sequence $\{(\s_i,j_i)\}_{i=1}^{N+2}$ one has that if there exists $M=M(N)$ such that
%	\[
%	\sum_{i=1}^{N+2} \s_i \omega_{k,j_i} =0 \,,\quad \forall k \leq M
%	\]
%	implies that $N$ is even and,  up to a permutation,
%	\[
%	\{(\s_i,j_i)\}_{i=1}^{N+2} =\{(+,\hat\jmath_1),\dots,(+,\hat\jmath_{N/2+1}), (-,\hat\jmath_1),\dots,(-,\hat\jmath_{N/2+1})\}
%	\]




\begin{prop}\label{formal-lin}Consider a Family of Hamiltonians satisfying H1 and H2.
	There exists $S\in  \cF^{\geq 1}$ such that for all $k\in \N_{\geq1}$ one has 
	\[
	e^{\{S,\cdot\}} H_k = L_k  + Z_k\,,\quad Z_k \in \cK^{\geq 2}\,.
	\]
\end{prop}
\begin{proof}
	We proceed by induction. Let us assume that, for some $n\geq 1$ there exists an $S_n$ such that
	for all $k$
	\[
	e^{\{S_n,\cdot\}} H_k=:H_{k,n}= L_k  + Z_{k,n}+ R_{k,n}\,,\quad Z_{k,n} \in \cK^{2\leq  \td < n}\,, \quad R_{k,n} \in \cR^{\geq n}\,.
	\]
	In order to make a normal form step we look for a formal change of variables $e^{\{F_n,\cdot\}}$, with $F_n\in \cR^{(n)}$ so that
	\[
	e^{\{F_n,\cdot\}} H_k = L_k  + Z_{k,n+1}+ R_{k,n+1 },,\quad Z_{k,n} \in \cK^{2\leq  \td < n+1}\,, \quad R_{k,n} \in \cR^{\geq n+1}
	\]
	By construction 
	\begin{align*}
		e^{\{F_n,\cdot\} }H_{k,n} &=L_k+ Z_{k,n} + R_{k,n} + \{F_n, L_k\} + \sum_{h=2}^\infty \frac{\ad_{F_n}^{h-1}}{h!} \{F_n,L_k \}
		+ \sum_{k=1}^\infty \frac{\ad_{F_n}^k}{k!}(Z_{k,n} + R_{k,n}) \\
		&= L_k+ Z_{k,n} + \Pi^\cK R_{k,n}  - \sum_{k=1}^\infty \frac{\ad_{F_n}^{k}}{(k+1)!}\Pi^{\cR}  P_{i}  + \sum_{k=1}^\infty \frac{\ad_{F_n}^k}{k!}(Z_{k,n} + R_{k,n})\,.
	\end{align*}
	So we may set
	\[
	Z_{k,n+1}:= Z_{k,n}+%\Pi^{\leq 2 i +2}
	\Pi^{(n)}\Pi^\cK R_{k,n} \,,\quad R_{k,n+1}= e^{\{F_n,\cdot\} }H_{k,n}- L_k - Z_{k,n+1}\,.
	\]
	In order to iterate our inductive step we need to show that we may fix $F_n$ so that for all $k$ 
	\[
	\Pi^{\td\leq  n} R_{k,n+1}= \{F_n,L_k\} + \Pi^\cR R_{k,n}^{(n)} =0\,.
	\]
	Passing to the Taylor coefficients this reads
	\begin{equation}
		\label{homotutte}
		\im \omega_k \cdot (\al-\bt)(F_n)_{\al,\bt} =(R_{k,n})_{\al,\bt} \,,\quad \forall \al\ne \bt\,:\quad |\al|+|\bt| =n+2\,.
	\end{equation}
	Finally the condition $\{H_k,K_j\}=0$ ensures that the equation above is trivial unless
	\begin{equation}
		\label{condk}
		{\al_j-\bt_j = \al_{-j} -\bt_{-j}}
	\end{equation}
	We claim that a  common solution to  the above infinite list of equations is given by
	\begin{equation}
		\label{questo}
		(F_n)_{\al,\bt} =\begin{cases}
			\im \dfrac{(R_{1,n})_{\al,\bt}}{\omega_1 \cdot (\al-\bt)}\,,\qquad \mbox{if}\;  \omega_1 \cdot (\al-\bt)\neq 0\\
			\im 	\dfrac{(R_{2,n})_{\al,\bt}}{\omega_2 \cdot (\al-\bt)}\,,\qquad \mbox{if}\;  \omega_1 \cdot (\al-\bt)=0 \;\mbox{and}\;   \omega_2 \cdot (\al-\bt)\neq 0\\
			\vdots
			\\
			\im 	\dfrac{(R_{k,n})_{\al,\bt}}{\omega_k \cdot (\al-\bt)}\,,\qquad \mbox{if}\;  \omega_1j\cdot (\al-\bt)=0 \; \forall j< k \;\mbox{and}\;   \omega_k \cdot (\al-\bt)\neq 0\\
			\vdots
		\end{cases}
	\end{equation}
	for all $(\al,\bt) \in \cM$ satisfyting \eqref{condk} and  $(F_n)_{\al,\bt}=0$ otherwise.
	To prove our claim, we use the fact that, by the involution hypothesis H1., we have 
	\[
	0= e^{\{S_n,\cdot \}} \{H_{k_1},H_{k_2}\}= \{H_{k_1,n},H_{k_2,n}\}=\{ L_{k_1}  + Z_{k_1,n}+ R_{k_1,n},L_{k_2}  + Z_{k_2,n}+ R_{k_2,n}\}
	\]
	so that, since $\{Z_{k_1,n},Z_{k_2,n}\}=0$ and $R_{k_i,n} \in \cF^{\geq n}$, one has 
	\[
	\Pi^{(n)} \{H_{k_1,n},H_{k_2,n}\} =\{L_{k_1}, R_{k_2,n}^{(n)}\} + \{R_{k_1,n}^{(n)}, L_{k_2}\}=0\,,
	\]
	passing to the Taylor coefficients we have
	\begin{equation}
		\label{commuto}
		\omega_{k_2} \cdot (\al-\bt) (R_{k_1,n})_{\al,\bt} = \omega_{k_1} \cdot (\al-\bt) (R_{k_2,n})_{\al,\bt} \,,\qquad \forall \; |\al|+|\bt| = n+2\,.
	\end{equation}
	Denoting $\ell= \al-\bt$, we recall that, by the independence assumption H2., there exists ${\mathbf k}_0 = {\mathbf k}_0(\al-\bt)$ such that
	\[
	\omega_{h} \cdot (\al-\bt) =0\quad  \forall h< {\mathbf k}_0\,,\qquad  \omega_{{\mathbf k}_0}\cdot (\al-\bt) \neq 0\,.
	\]
	By \eqref{commuto}, for all $k$ and all $ |\al|+|\bt| = n+2$ we have
	\[
	(R_{k,n})_{\al,\bt} =% \big( \omega_{{\mathbf k}_0}\cdot (\al-\bt)\big)^{-1}
	\frac{\omega_{k} \cdot (\al-\bt)}{\omega_{{\mathbf k}_0}\cdot (\al-\bt)} (R_{{\mathbf k}_0,n})_{\al,\bt} \,,\qquad \forall \; |\al|+|\bt| = n+2\,,
	\]
	which shows that $F_n$ defined in \eqref{questo} satisfies \eqref{homotutte}.
	Finally by Lemma \ref{compo} we can define $S_{n+1}\in \cF^{\geq 1}$ so that
	\[
	e^{\{S_{n+1},\cdot\}}= e^{\{F,\cdot \} }e^{\{S_n,\cdot\}}\,,
	\]
	passing to the limit we obtain the desired result.
\end{proof}
\subsection{Application to the Kirchhoff-Pohozaev hierarchy}
In order to prove our main result Theorem \ref{formal-birk},  it is convenient to pass  to complex notation, as seen in Section \ref{complessi}. For the expressions of the constants of motion we refer to  \eqref{costanticompl}.
By construction all the building blocks $D_i,B_i,Q_i$ are  a formal quadratic polynomials which preserve momentum. Thus they all belong to $\cF^{(0)}$ and,  by Lemma \ref{anello} , so does $A= (1+D_1)^{-1}$ as well as the $E_i$ and $F_i$. Again by Lemma \ref{anello} this ensures that the $\J_i\in \cF^{(0)}$ and moreover 
\[
\J_k = 2\sum_{j\in \Z_0} |j|^{2k-1}|z_j|^2 +\J_k^{\geq 2}\,,\qquad   \cI_k = |j|(|z_j|^2+|z_{-j|^2} )+\cI_k^{\geq 2}\,,\qquad  \J_k^{\geq 2},\cI_k^{\geq 2}\in \cF^{\geq 2}\,.
\] 
We claim that both the families  $\J_k$ and $\cI_k$ satisfy Assumptions H1, H2. The involution assumption is a direct consequence of Proposition \ref{involuz} and  Lemma \ref{prop:comm_gamma2}. Regarding assumption H2, for the $\cI_|j|^2$  this is obvious, since the $(|z_j|^2+|z_{-j}|^2 ,|z_j|^2-|z_{-j}|^2)_{j\in \N_{\geq 1}}$ are clearly independent.
\\
Regarding the $\I_k$, let us show that, setting $L_k:= \sum_{j\in \Z_0} |j|^{2k-1}|z_j|^2 $, and hence $\omega_k = |j|^{2k+1}$, Assumption H2 holds.

This is proved for instance in \cite{BambusiGrebert2006TameModulus}.
In fact, assume by contradiction that H2 does not hold. Since $\ell\in \Z^\Z_f$ has finite support and $\ell_j= \ell_{-j}$, we may assume, without loss of generality, that $\ell=(\ell_{j_i})_{i=1,\dots,M}$ for some $M$ with all the $\ell_{j_i}\ne 0$ and distinct $j_i>0$.  Let us then consider the system of $M$  linear equations
\[
\begin{cases}
	&\sum_{i=1}^M \omega_{1,j_i} x_i =0 \\&\vdots\\
	&\sum_{i=1}^M \omega_{M,j_i} x_i =0 \,.
\end{cases}
\] 
If H2 were false, there would exist a non trivial solution to the system above, which in turm means that 
\[
\det(\Omega)=0 \,,\qquad \Omega= \begin{pmatrix}
	j_1 &j_2 &\dots & j_M \\
	j_1^3 &j_2^3 &\dots & j_M^3\\
	\vdots &\vdots &\vdots&\vdots\\
	j_1^{2M-1} &j_2^{2M-1} &\dots & j_M^{2M-1}
\end{pmatrix}
\]
but this cannot hold true, since this is a Vandermonde matrix we have 
\[
\det(\Omega) = \left( \prod_{i=1}^M j_i \right) \prod_{1 \leq  k < i \leq M} (j_i^2 - j_k^2)
\]
and the result follows. In both cases we apply Proposition \ref{formal-lin}, obtaining in principle two formal normalizing changes of variables. Denote by $\Psi$ the change of variables which normalizes the $\cI$ hierarchy. In order to show that it normalizes also the $\J_k$ we only need to remark that the identities  \eqref{calIJ} holds at the level of formal power series.
\bibliographystyle{plain}
\bibliography{BibliografiaTot}
\end{document}
